\documentclass[11pt,reqno]{amsart}

\usepackage[T1]{fontenc}
\usepackage[utf8]{inputenc}
\usepackage{lmodern}
\usepackage{microtype}
\usepackage{amsmath,amssymb,mathtools}
\usepackage{booktabs}
\usepackage{enumitem}
\usepackage[margin=1.05in]{geometry}
\usepackage{xcolor}
\usepackage[colorlinks=true,linkcolor=blue!55!black,
 citecolor=blue!55!black,urlcolor=blue!65!black]{hyperref}
\hypersetup{pdftitle={Trace Sparsity and Arithmeticity for Compact Hyperbolic Surfaces},pdfauthor={}}
\ifdefined\pdftrailerid
  \pdftrailerid{}
\fi

\newtheorem{theorem}{Theorem}[section]
\newtheorem{proposition}[theorem]{Proposition}
\newtheorem{lemma}[theorem]{Lemma}
\newtheorem{corollary}[theorem]{Corollary}
\newtheorem{criterion}[theorem]{Criterion}
\theoremstyle{definition}
\newtheorem{definition}[theorem]{Definition}
\newtheorem{convention}[theorem]{Convention}
\theoremstyle{remark}
\newtheorem{remark}[theorem]{Remark}
\newtheorem{warning}[theorem]{Warning}

\newcommand{\R}{\mathbb R}
\newcommand{\C}{\mathbb C}
\newcommand{\Q}{\mathbb Q}
\newcommand{\Z}{\mathbb Z}
\newcommand{\F}{\mathbb F}
\newcommand{\cO}{\mathcal O}
\newcommand{\cB}{\mathcal B}
\newcommand{\cG}{\mathcal G}
\newcommand{\cJ}{\mathcal J}
\newcommand{\Tr}{\operatorname{tr}}
\newcommand{\GL}{\operatorname{GL}}
\newcommand{\SL}{\operatorname{SL}}
\newcommand{\PSL}{\operatorname{PSL}}
\newcommand{\Spec}{\operatorname{Spec}}
\newcommand{\Res}{\operatorname{Res}}
\newcommand{\Gal}{\operatorname{Gal}}
\newcommand{\Norm}{\operatorname{N}}

\title[Trace sparsity and arithmeticity]
 {Trace Sparsity and Arithmeticity for Compact Hyperbolic Surfaces}
\author{}
\date{Working paper, 26 September 2026}

\begin{document}

\begin{abstract}
Let \(\Gamma<\PSL_2(\R)\) be a torsion-free cocompact lattice, and let
\(T(\Gamma)\) be the set of traces of the full inverse image of \(\Gamma\)
in \(\SL_2(\R)\).  We prove arithmeticity already at the critical
trace-growth exponent:
\[
 \limsup_{X\to\infty}
 \frac{\log\#\bigl(T(\Gamma)\cap[-X,X]\bigr)}{\log X}\le1
 \quad\Longrightarrow\quad
 \Gamma\text{ is arithmetic}.
\]
Equivalently, the hypothesis is the family of bounds
\(\#(T(\Gamma)\cap[-X,X])=O_{\Gamma,\eta}(X^{1+\eta})\) for every
\(\eta>0\).  In particular, linear growth and bounded clustering of the
signed trace set force arithmeticity.  These consequences are the compact
trace formulations of the conjectures of Schmutz and Sarnak.

The main estimate compares translation length at the original Fuchsian
character with arithmetic expansion at an algebraic specialization.
Let \(x\) be the Fuchsian character of \(\Gamma\), let \(y\) be an algebraic
Fuchsian character in the rational Zariski closure of \(x\), and let \(k_y\)
be the trace field of the subgroup generated by squares.  For every element
\(b\) of that subgroup we prove
\[
 \cJ_{k_y}(\Tr_y b)\le \frac{\ell_x(b)}2,
\]
where \(\cJ_k\) is the sum of the local expansion exponents of the associated
Chebyshev recurrence.  When the original character is algebraic, the
distinguished real summand of \(\cJ_{k_x}(\Tr_x b)\) already equals
\(\ell_x(b)/2\).  Equality therefore holds, all other local exponents
vanish, and Takeuchi's criterion gives arithmeticity.

To prove this estimate, we fix an index in the Chebyshev recurrence and
vary the group element.  Roth's theorem bounds from below the reduced
divisor of two consecutive recurrence terms.  Expansion in congruence
quotients modulo \(p^2\) allows us to choose elements for which no prime
occurs in that divisor to an unbounded power.  We obtain the corresponding
upper bound by combining three counts: a uniform bound for fibers of two
trace functions, a congruence count for moduli with bounded prime
exponents, and an injective map from length-two paths in a bipartite trace
graph.  Algebraic specialization and Thurston's asymmetric
length-spectrum metric then show that the original Fuchsian character
cannot be transcendental.  Section~\ref{sec:external-inputs} states the
analytic and arithmetic results from the literature, including all
uniformity needed in these counts.
\end{abstract}

\maketitle
\enlargethispage{2pt}
\tableofcontents

\section{Introduction}

Let \(S\) be a closed orientable hyperbolic surface and let
\(\Gamma<\PSL_2(\R)\) be the corresponding torsion-free cocompact lattice.
Write \(\widetilde\Gamma\) for its full inverse image in \(\SL_2(\R)\), and
define the \emph{signed trace set}
\[
 T(\Gamma)=\{\Tr g:g\in\widetilde\Gamma\}.
\]
Thus \(T(\Gamma)=-T(\Gamma)\); it is a set of values, and repeated traces are
not counted with multiplicity.  Put
\[
 N_\Gamma(X)=\#\bigl(T(\Gamma)\cap[-X,X]\bigr)
 \quad (X\ge2)
\]
and define its \emph{upper trace-growth exponent} by
\begin{equation}\label{eq:upper-trace-exponent}
 \overline\delta_T(\Gamma)
 =\limsup_{X\to\infty}\frac{\log N_\Gamma(X)}{\log X}.
\end{equation}
We say that the trace set satisfies the \emph{critical trace-growth bound}
if \(\overline\delta_T(\Gamma)\le1\).  This is equivalent to
\begin{equation}\label{eq:critical-trace-growth}
 N_\Gamma(X)=O_{\Gamma,\eta}(X^{1+\eta})
 \qquad\text{for every }\eta>0,
\end{equation}
or, more briefly, to the upper bound \(N_\Gamma(X)\le X^{1+o(1)}\).
We say that \(T(\Gamma)\) has
\emph{linear growth} if
\begin{equation}\label{eq:linear-growth}
 \#\bigl(T(\Gamma)\cap[-X,X]\bigr)=O_\Gamma(X)
 \qquad (X\longrightarrow\infty).
\end{equation}
It has \emph{bounded clustering} if
\begin{equation}\label{eq:bounded-clustering}
 \sup_{u\in\R}\#\bigl(T(\Gamma)\cap[u-1,u+1]\bigr)<\infty.
\end{equation}
Linear growth implies the critical trace-growth bound.  The choice of radius
in \eqref{eq:bounded-clustering} is immaterial: an
interval of one fixed positive radius can be covered by finitely many
intervals of any other fixed positive radius.  Covering \([-X,X]\) by
\(O(X)\) intervals of radius one then shows that bounded clustering implies
linear growth.

If \(g\in\widetilde\Gamma\) projects to a nonidentity element, its translation
length satisfies
\begin{equation}\label{eq:trace-length}
 |\Tr g|=2\cosh\frac{\ell(g)}2.
\end{equation}
Consequently, \(N_\Gamma(X)\) differs by at most a factor of two from the
number of distinct translation lengths at most
\(2\operatorname{arcosh}(X/2)=2\log X+O(1)\); the factor two accounts only
for the two trace signs.  This observation converts the trace statements used here
into the usual length formulations in \cite{Sarnak,Schmutz}.

For context, Luo and Sarnak proved that arithmetic cofinite Fuchsian groups
have bounded clustering \cite{LuoSarnak}.  Geninska and Leuzinger proved
the converse under bounded clustering for nonuniform lattices
\cite{GeninskaLeuzinger}, and Hao subsequently proved the stronger
linear-growth converse in the nonuniform case \cite{HaoTraceSets}.  For
closed surfaces of genus at least six, Hao also proved a superlinear
generic lower bound outside a countable union of positive-codimension
algebraic subsets of Teichm\"uller space \cite{HaoLengthSets}.  The compact
case for all hyperbolic metrics addressed here was not covered by those
results.

We recall the arithmetic terminology used in the conclusion.  A quaternion
algebra \(B\) over a number field \(k\) is a four-dimensional central simple
\(k\)-algebra.  It has a reduced norm
\(\operatorname{Nrd}:B\to k\), which becomes the determinant after any
field extension that identifies \(B\) with \(M_2\).  In the Fuchsian case,
\(B/k\) is called \emph{admissible} when \(k\) is totally real, \(B\) is
split at one distinguished real place, and \(B\) is ramified at every other
real place.  An \emph{order} \(\mathcal O\subset B\) is a subring containing
\(1\) which is also a full \(\cO_k\)-lattice.  Its norm-one group is
\[
 \mathcal O^1=\{u\in\mathcal O:\operatorname{Nrd}(u)=1\}.
\]
Using the distinguished splitting \(B\otimes_{k,v_0}\R\simeq M_2(\R)\),
the projective image of \(\mathcal O^1\) is a lattice in
\(\PSL_2(\R)\).  A Fuchsian group is \emph{arithmetic} if it is
commensurable with such a projective norm-one group.  Two groups are
commensurable when their intersection has finite index in each.  A group is
\emph{derived from a quaternion algebra} when, after conjugation, it is
contained with finite index in the projective norm-one group of an order.
These formulations define the same commensurability class.

Our main theorem is the following.

\begin{theorem}[trace sparsity forces arithmeticity]
\label{thm:main-arithmeticity}
Let \(\Gamma<\PSL_2(\R)\) be torsion-free and cocompact.  If its signed
trace set satisfies \(\overline\delta_T(\Gamma)\le1\), then \(\Gamma\) is
arithmetic.
\end{theorem}

\begin{corollary}[compact Schmutz and Sarnak conjectures and an orbifold extension]
\label{cor:schmutz-sarnak}
Let \(\Gamma<\PSL_2(\R)\) be any cocompact Fuchsian lattice, possibly with
torsion.  The critical trace-growth bound forces \(\Gamma\) to be arithmetic.
In particular, linear growth or bounded clustering of the signed trace set
forces arithmeticity.  Thus the two compact-surface conjectures hold, and
the same conclusions extend to compact orientable hyperbolic orbifolds.
\end{corollary}

Theorem~\ref{thm:main-arithmeticity} is the torsion-free part of
Corollary~\ref{cor:schmutz-sarnak}.  Selberg's lemma and invariance of
arithmeticity under commensurability give the extension to lattices with
torsion.  Linear growth implies the critical bound, while the covering
argument following \eqref{eq:bounded-clustering} shows that bounded
clustering implies linear growth.  The main step is a local inequality
which bounds all arithmetic expansion rates at an algebraic Fuchsian
character by the geometric length at the original character.

Let \(\Pi=\pi_1(S)\), and let
\[
 \Pi^{(2)}=\langle\gamma^2:\gamma\in\Pi\rangle.
\]
The quotient \(\Pi/\Pi^{(2)}\) is the finite elementary abelian group
\(H_1(S,\F_2)\), so \(\Pi^{(2)}\) has finite index.  A lift of a
\(\PSL_2\)-representation to \(\SL_2\) is not unique, but two lifts differ
by a character \(\Pi\to\{\pm1\}\).  Such a character is trivial on
\(\Pi^{(2)}\).  Hence traces on \(\Pi^{(2)}\), and the field generated by
them, do not depend on the chosen lift.
For a marked Fuchsian group \(\Gamma\), we similarly write
\[
 \Gamma^{(2)}=\langle g^2:g\in\Gamma\rangle;
\]
it is the projective image of \(\Pi^{(2)}\).  If
\(\widetilde\Gamma<\SL_2(\R)\) is the full inverse image, then
\[
 k_\Gamma=\Q\bigl((\Tr \widetilde g)^2:
                         \widetilde g\in\widetilde\Gamma\bigr)
\]
denotes the invariant trace field.

For a number field \(k\) and \(t\in k\), Section~\ref{sec:local-expansion}
defines a nonnegative quantity \(\cJ_k(t)\).  When a distinguished real
embedding sends \(t\) to the trace of a hyperbolic element of translation
length \(L\), that one embedding contributes exactly \(L/2\) to
\(\cJ_k(t)\).

\begin{theorem}[local expansion for algebraic points in the rational closure]
\label{thm:local-expansion}
Let \(x\) be a faithful cocompact Fuchsian character of the closed surface
group \(\Pi\), and assume that its signed trace set satisfies the critical
trace-growth bound.  Let \(\mathcal X\) be the affine \(\Q\)-variety of
liftable \(\PSL_2\)-characters described in
Section~\ref{sec:characters}, and let \(Z_x\subset\mathcal X\) be the
rational Zariski closure of \(x\).  Let
\(y\in Z_x(\overline\Q)\) be an algebraic faithful cocompact
Fuchsian character, and put
\[
 k_y=\Q\bigl(\Tr_y h:h\in\Pi^{(2)}\bigr).
\]
Then, for every \(b\in\Pi^{(2)}\),
\begin{equation}\label{eq:main-local}
 \cJ_{k_y}(\Tr_y b)\le\frac{\ell_x(b)}2.
\end{equation}
For the identity element both sides are zero.
\end{theorem}

\begin{corollary}[rigidity among algebraic Fuchsian specializations]
\label{cor:specialization-rigidity}
Under the hypotheses of Theorem~\ref{thm:local-expansion}, assume in addition
that \(y\) lies in the same Fuchsian component as \(x\).  This holds, in
particular, for the nearby specializations supplied by
Proposition~\ref{prop:bounded-degree-specialization}.  Then the two marked
Fuchsian characters are equal: \(y=x\).
\end{corollary}

\begin{proof}
For every \(\gamma\in\Pi\), apply Theorem~\ref{thm:local-expansion} to
\(b=\gamma^2\in\Pi^{(2)}\).  The distinguished real place of \(k_y\)
contributes \(\ell_y(b)/2\) to the nonnegative sum
\(\cJ_{k_y}(\Tr_y b)\).  Hence
\[
 \ell_y(\gamma)
 =\frac{\ell_y(\gamma^2)}2
 \le \cJ_{k_y}(\Tr_y(\gamma^2))
 \le\frac{\ell_x(\gamma^2)}2
 =\ell_x(\gamma).
\]
Thurston's minimal-stretch theorem identifies the logarithm of the least
Lipschitz constant in the marked homotopy class with
\[
 d_{\rm Th}(x,y)
 =\log\sup_{\gamma\ne1}\frac{\ell_y(\gamma)}{\ell_x(\gamma)}.
\]
It defines a separating asymmetric metric on Teichm\"uller space
\cite{ThurstonStretch}.  The displayed domination therefore gives
\(d_{\rm Th}(x,y)\le0\).  Since this metric is nonnegative, it is zero, and
its separation property gives \(x=y\).
\end{proof}

Here and below an element of \(\Pi^{(2)}\) and its lift to \(\SL_2\) are
denoted by the same letter when the choice of sign is irrelevant.  Since
\(\cJ_k(-t)=\cJ_k(t)\), inequality \eqref{eq:main-local} is independent of
that sign.

We next fix notation used in the proof.  For a
number field \(k\) and a finite set \(S\) of its places, let
\begin{equation}\label{eq:S-integer-definition}
 \cO_{k,S}=\{a\in k: |a|_v\le 1
       \text{ for every finite place }v\notin S\}
\end{equation}
be the ring of \(S\)-integers.  Thus archimedean places may be included in
\(S\), but do not affect the ring.  A nonzero integral ideal of
\(\cO_{k,S}\) has a unique factorization
\(\mathfrak a=\prod_{\mathfrak p\notin S}\mathfrak p^{e_{\mathfrak p}}\),
and its absolute norm is
\begin{equation}\label{eq:ideal-norm-definition}
 \Norm\mathfrak a
 =\#(\cO_{k,S}/\mathfrak a)
 =\prod_{\mathfrak p\notin S}
   \#(\cO_k/\mathfrak p)^{e_{\mathfrak p}}.
\end{equation}
In particular, places in \(S\) do not contribute to the norm of a principal
\(S\)-ideal.

Choose a basepoint \(o\in\mathbb H^2\).  The notation
\(d_x(o,go)\) means hyperbolic displacement under the representation with
character \(x\), and \(\ell_x(g)=\inf_z d_x(z,gz)\) is its translation
length.  Word length, denoted by \(|g|\), is always taken with respect to a
finite symmetric generating set fixed before the relevant argument.  In the
main construction we use the image under \(\rho_y\) of one fixed generating
set of \(\Pi^{(2)}\).  Replacing this generating set changes only constants,
because any two word metrics on a finitely generated group are
quasi-isometric.

Subscripts on asymptotic notation specify the data on which the constants
may depend.
For example, \(r_m=o_{\mathcal D}(m)\) means
\(r_m/m\to0\) while the data \(\mathcal D\) remain fixed, and
\(O_{\mathcal D}(1)\) denotes a quantity bounded in terms of those data.
An unsubscripted \(O\)- or \(o\)-term inside a proof is uniform after all
objects fixed earlier in that proof have been fixed.  Whenever auxiliary
data vary in a later limiting operation, we state the order of limits
explicitly.

Here is the argument for Theorem~\ref{thm:local-expansion} in outline.
For
\[
 q_0(T)=0,\qquad q_1(T)=1,\qquad
 q_{n+1}(T)=Tq_n(T)-q_{n-1}(T),
\]
fix \(N\ge3\) and put
\[
 F_N(T)=q_{N-1}(T)q_N(T).
\]
Starting from \(b\), we choose
\(a_m=b^m h_m\), where \(h_m\) has word length at most \(\theta m\) in the
fixed generating set.  If
\(s_m=\Tr_y(a_m)\), this choice has three properties:
\begin{align*}
 \cJ_{k_y}(s_m)
   &\ge m\cJ_{k_y}(\Tr_y b)-O(\theta m)-O(1),\\
 d_x(o,a_mo)
   &\le m\ell_x(b)+O(\theta m)+O(1),\\
 F_N(s_m)\cO_{k_y,S}
   &\text{ has uniformly bounded prime exponents.}
\end{align*}
Roth's theorem then gives the lower estimate
\begin{equation}\label{eq:intro-roth}
 \log\Norm\bigl(F_N(s_m)\cO_{k_y,S}\bigr)
 \ge(2N-4-\epsilon)\cJ_{k_y}(s_m)-O(1).
\end{equation}
The trace-graph and congruence argument gives the upper estimate
\begin{equation}\label{eq:intro-wedge}
 \log\Norm\bigl(F_N(s_m)\cO_{k_y,S}\bigr)
 \le\frac{2N-1}{2}d_x(o,a_mo)+o(m).
\end{equation}
After division by \(m\), we let \(m\to\infty\), then
\(\theta,\epsilon\to0\), and finally \(N\to\infty\).  Since
\[
 \frac{(2N-1)/2}{2N-4-\epsilon}\longrightarrow\frac12,
\]
this proves \eqref{eq:main-local}.

The limits must be taken in this order.  The finite set of excluded
places, the exponent bound in the sieve, and the ball radius used for
congruence counting may depend on \(N,\theta,\epsilon,b,y\), but they are
fixed before \(m\) tends to infinity.

The following sections prove these estimates.  Section~\ref{sec:external-inputs}
states the results from the literature in the forms used here.  The
superapproximation needed to select \(a_m\) concerns only squares of
primes; the later congruence count uses composite moduli whose prime
exponents are bounded by one fixed constant.  No expansion theorem for
arbitrary prime powers or arbitrary ideals is assumed.

\section{Fuchsian characters and rational closure}
\label{sec:characters}

\subsection{Characters, lifts, and trace functions}

Fix a presentation of the closed surface group \(\Pi\).  Let
\(\mathcal X\) be the affine \(\Q\)-variety of liftable
\(\PSL_2\)-characters obtained from the \(\SL_2\)-character variety by
quotienting by the finite group of sign twists
\(\operatorname{Hom}(\Pi,\{\pm1\})\).  Thus \(\Q[\mathcal X]\) is the
invariant subring of the \(\SL_2\)-character algebra.  It is generated by
the adjoint characters, or equivalently by the functions
\(\tau_\gamma=\Tr^2(\gamma)\).  Finite generation of
matrix invariants and the trace algebra permits a finite collection of
squared-trace generators \cite{Procesi,Horowitz}.  For the corresponding
description of the \(\PSL_2\)-character algebra, see
\cite[Proposition~2.2 and Theorem~1.3]{HeusenerPorti}.
Those results are stated over \(\C\): Proposition~2.2 gives generation by
the functions \(\tau_\gamma\), and Theorem~1.3 identifies the closed points
with squared-trace characters.  Faithfully flat descent along
\(\C/\Q\) gives the same generation statement for the chosen
\(\Q\)-model: a quotient of finitely generated \(\Q\)-algebras is an
isomorphism if it becomes an isomorphism after tensoring with \(\C\).
Since the coordinate ring is finitely generated, finitely many words
suffice.
After an orientation has been fixed, the real Fuchsian component is naturally
identified with Teichm\"uller space.  Goldman's description by the extremal
Euler-class components shows that this is a connected component of the real
representation space; after quotienting by conjugation, it is therefore open
in the smooth real locus of \(\mathcal X\)
\cite[Corollary~C and Section~10.4]{GoldmanComponents}.
For \(\gamma\in\Pi\), the function
\[
 \tau_\gamma(z)=\Tr_z(\gamma)^2
\]
is a well-defined regular function on \(\mathcal X\), independent of a
choice of lift.  If \(h\in\Pi^{(2)}\), then \(\Tr_z(h)\) itself is
independent of the lift and is regular on \(\mathcal X\).
Alternatively, one may work on a fixed finite cover that parameterizes
\(\SL_2\)-lifts.  Since the degree of this cover is fixed, passing to it
changes each cardinality estimate below by at most a fixed factor.  The
finite generation statement we need is the following.

\begin{lemma}[finite squared-trace coordinates and the invariant trace field]
\label{lem:finite-trace-coordinates}
There are words \(\gamma_1,\ldots,\gamma_r\in\Pi\) such that the functions
\(\tau_{\gamma_i}\) generate \(\Q[\mathcal X]\) as a \(\Q\)-algebra.  For
every nonelementary character \(y\),
\begin{equation}\label{eq:invariant-trace-field-identity}
 \begin{split}
 k_y
 &=\Q\bigl(\Tr_y h:h\in\Pi^{(2)}\bigr)\\
 &=\Q\bigl(\Tr_y(\gamma)^2:\gamma\in\Pi\bigr)
  =\Q\bigl(\tau_{\gamma_1}(y),\ldots,
             \tau_{\gamma_r}(y)\bigr).
 \end{split}
\end{equation}
In particular, if \(y\) is an algebraic point with residue field \(\Q(y)\),
then
\begin{equation}\label{eq:trace-field-degree-bound}
 k_y\subseteq\Q(y),\qquad
 [k_y:\Q]\le[\Q(y):\Q].
\end{equation}
\end{lemma}

\begin{proof}
By the choice of the affine model, its coordinate ring is generated by
squared-trace functions.  Since that ring is finitely generated, finitely
many of them suffice; call them \(\tau_{\gamma_i}\).  The first equality in
\eqref{eq:invariant-trace-field-identity} is the definition of the full
invariant trace field.  The second is the invariant-trace-field identity for
a nonelementary subgroup of \(\PSL_2\); see
\cite[Lemma~3.5.6]{MaclachlanReid}.  One inclusion is also visible from
\[
 \Tr_y(\gamma^2)=\Tr_y(\gamma)^2-2,
 \qquad \gamma^2\in\Pi^{(2)}.
\]
The finite set \(\tau_{\gamma_i}\) generates every squared-trace function
as an element of the coordinate ring, which gives the last equality.
Evaluation of regular functions at an algebraic point takes values in its
residue field, proving \eqref{eq:trace-field-degree-bound}.
\end{proof}

\subsection{Rational closure and persistence of trace collisions}

Let \(A=\Q[\mathcal X]\).  For \(x\in\mathcal X(\R)\), define
\begin{equation}\label{eq:rational-ideal}
 I(x)=\{f\in A:f(x)=0\},\qquad
 Z_x=V(I(x))\subset\mathcal X.
\end{equation}
We call \(Z_x\) the \emph{rational Zariski closure} of \(x\).  The
terminology refers to the fact that it is the smallest closed subvariety
defined over \(\Q\) that contains \(x\).

\begin{lemma}[trace equalities persist on the rational closure]
\label{lem:collision-persistence}
If \(y\in Z_x\) and
\(\tau_\alpha(x)=\tau_\beta(x)\), then
\(\tau_\alpha(y)=\tau_\beta(y)\).  More generally, every polynomial
relation over \(\Q\) among trace functions that holds at \(x\) holds at
every point of \(Z_x\).
\end{lemma}

\begin{proof}
The difference \(\tau_\alpha-\tau_\beta\) belongs to \(I(x)\), and every
element of \(I(x)\) vanishes on \(V(I(x))\).  The same argument applies to
an arbitrary polynomial relation.
\end{proof}

The use of squared traces in this lemma is important.  At a second point
\(y\), a signed equality may acquire either sign, but a fiber of squared
traces splits into at most two signed fibers.  Thus passing from squared
traces at \(x\) to signed traces at \(y\) increases the number of values by
at most a factor of two.

\begin{proposition}[algebraic specialization of bounded degree]
\label{prop:bounded-degree-specialization}
Let \(x\) be a Fuchsian point.  If \(\dim Z_x>0\), there are a positive
integer \(D\) and distinct algebraic Fuchsian points
\(y_j\in Z_x(\overline\Q)\cap\mathcal X(\R)\), all in the same Fuchsian
component as \(x\), such that
\begin{equation}\label{eq:specialization-convergence}
 y_j\longrightarrow x,\qquad [\Q(y_j):\Q]\le D.
\end{equation}
For every finite set \(E\subset\Pi\),
\begin{equation}\label{eq:partition-coarsening}
 \#\{\tau_\gamma(y_j):\gamma\in E\}
 \le \#\{\tau_\gamma(x):\gamma\in E\}.
\end{equation}
If the signed trace set at \(x\) satisfies the critical trace-growth bound and
\[
 \cB_R^x=\{\gamma\in\Pi:d_x(o,\gamma o)\le R\},
\]
then there is a function \(\omega_x(R)=o_x(R)\), independent of \(j\), such
that
\begin{equation}\label{eq:specialized-trace-support}
 \#\{\Tr_{y_j}\gamma:\gamma\in\cB_R^x\}
 \le \exp\bigl(R/2+\omega_x(R)\bigr).
\end{equation}
The same assertion holds after replacing \(\Pi\) by any fixed finite-index
subgroup.
\end{proposition}

\begin{proof}
The ideal \(I(x)\) is the kernel of the evaluation homomorphism
\(A\to\R\), hence it is prime and \(Z_x\) is irreducible.  Moreover, \(x\)
does not lie on a proper closed subvariety of \(Z_x\) defined over \(\Q\):
an additional equation defining such a subvariety would vanish at \(x\) and
would therefore already belong to \(I(x)\).  Since the ground field has
characteristic zero, the singular locus is a proper \(\Q\)-closed subset.
Thus \(x\) is a smooth point of \(Z_x\).

Put \(d=\dim Z_x\).  Choose a rational linear Noether normalization
\[
 \pi:Z_x\longrightarrow\mathbb A^d
\]
whose differential is invertible at \(x\).  Such a choice exists because
the admissible linear projections form a nonempty Zariski-open set, and
rational points are dense in that parameter space.  After removing the
branch and nonflat loci, \(\pi\) is finite and \'etale in a neighborhood of
\(x\).  The real inverse-function theorem gives a local real analytic
inverse branch through \(x\).  Rational points are dense in
\(\mathbb A^d(\R)\), so one may choose distinct rational points approaching
\(\pi(x)\) within the image of that branch and lift them along the branch.
Each lifted point lies in a fiber of the finite morphism over a rational
point.  It is therefore algebraic of degree at most \(\deg\pi\), and the
lifted points converge to \(x\).  Since the Fuchsian component containing
\(x\) is open, the branch may be restricted so that every \(y_j\) is
faithful, cocompact, and lies in the same Fuchsian component as \(x\).

Lemma~\ref{lem:collision-persistence} says that two elements of \(E\) with
the same squared trace at \(x\) still have the same squared trace at
\(y_j\).  Thus the partition of \(E\) by squared trace at \(y_j\) is
coarser than the corresponding partition at \(x\), which proves
\eqref{eq:partition-coarsening}.  If
\(\gamma\in\cB_R^x\), then
\[
 |\Tr_x\gamma|\le 2\cosh(R/2).
\]
The critical trace-growth bound therefore gives
\[
 \#\{\Tr_x\gamma:\gamma\in\cB_R^x\}
 \le \exp\bigl(R/2+o_x(R)\bigr).
\]
Choose one function \(\omega_x(R)=o_x(R)\) that absorbs this error and the
factor two introduced when signs are restored.  The same estimate holds
for squared traces.
At \(y_j\), collision persistence shows that no more squared-trace values
occur, and each squared trace has at most two signed square roots.  This
proves \eqref{eq:specialized-trace-support} with a bound independent of
\(j\).  Passing to a fixed finite-index subgroup only restricts the set
being counted.
\end{proof}

If \(\dim Z_x=0\), then the residue field of the point is finite over
\(\Q\), so \(x\) is algebraic.  Proposition
\ref{prop:bounded-degree-specialization} is needed only when \(x\) is
transcendental.

\subsection{The trace field and its algebraic group}

Let \(y\) be an algebraic faithful Fuchsian character and let
\(\rho_y\) be one fixed \(\SL_2\)-lift.  Let
\(\Delta=\rho_y(\Pi^{(2)})<\SL_2(\R)\), and write
\(\overline\Delta<\PSL_2(\R)\) for its projective image.  Put
\[
 k=k_y=\Q(\Tr\Delta).
\]
The two groups play different roles in the congruence argument.  The linear
group \(\Delta\) carries the signed trace, the integral model, and the
reduction maps.  The projective group \(\overline\Delta\) is the
cocompact Fuchsian lattice whose finite-index subgroups define surface
covers.  A metric ball consists of abstract elements of \(\Pi^{(2)}\):
their displacement is measured through the projective representation at
\(x\), while their traces are evaluated through \(\rho_y\).
Because the projective representation at \(y\) is faithful,
\(\rho_y|_{\Pi^{(2)}}\) is injective, and we identify \(\Delta\) with
\(\Pi^{(2)}\) only as an abstract group.  Thus, for \(A\in\Delta\), its
\(x\)-displacement and word length mean those of the unique abstract element
\(\rho_y^{-1}(A)\).
Faithfulness of the projective representation, together with
torsion-freeness, implies
\(\Delta\cap\{\pm I\}=\{I\}\).  Passing instead to a full inverse image
adds at most the central element \(-I\) and changes all quotient and
cardinality estimates by a factor at most two.

\begin{lemma}[algebraic matrix realization and the associated quaternion algebra]
\label{lem:algebraic-realization}
After conjugation in \(\SL_2(\C)\), every matrix in \(\Delta\) has
algebraic entries.  Moreover, the \(k\)-span of \(\Delta\) is a
four-dimensional central simple \(k\)-algebra.
\end{lemma}

\begin{proof}
Every trace \(\Tr_y h\), for \(h\in\Pi^{(2)}\), is the value at the
algebraic point \(y\) of a regular function defined over \(\Q\), and is
therefore algebraic.  Choose \(A,C\in\Delta\) such that \(A\) is
hyperbolic and \(C\) preserves neither eigenline of \(A\).  Such a pair
exists because \(\Delta\) is nonelementary.  Conjugate \(A\) to
\[
 A=\begin{pmatrix}\lambda&0\\0&\lambda^{-1}\end{pmatrix}.
\]
The eigenvalue \(\lambda\) is algebraic because it satisfies
\(X^2-(\Tr A)X+1=0\).  Write
\(C=\left(\begin{smallmatrix}p&q\\r&s\end{smallmatrix}\right)\).
The two algebraic numbers \(\Tr C\) and \(\Tr(AC)\) determine \(p,s\).
Both \(q\) and \(r\) are nonzero, and \(qr=ps-1\) is algebraic.
Conjugation by a diagonal matrix, which leaves \(A\) fixed, may therefore
normalize \(q\) to \(\operatorname{sgn}(q)\): explicitly, use the real
matrix
\(\operatorname{diag}(|q|^{-1/2},|q|^{1/2})\).
The other off-diagonal entry then becomes
\(\operatorname{sgn}(q)\,qr\), and is therefore algebraic.  Thus this
real conjugation makes both \(A\) and \(C\) algebraic while preserving
the distinguished embedding into \(M_2(\mathbb R)\).

For \(H=\left(\begin{smallmatrix}x&y\\z&w\end{smallmatrix}\right)\in
\Delta\), the traces of \(H\) and \(AH\) determine \(x,w\).  After those
entries are known, the traces of \(CH\) and \(ACH\) give the two equations
\[
 qz+ry\in\overline\Q,\qquad
 \lambda qz+\lambda^{-1}ry\in\overline\Q.
\]
Their determinant is \(qr(\lambda^{-1}-\lambda)\ne0\), so \(y,z\) are
algebraic.  This proves the first assertion.

The four matrices \(I,A,C,AC\) form a basis of \(M_2(\C)\): the two
diagonal entries are separated by \(I,A\), and the two off-diagonal entries
are separated by \(C,AC\).  The bilinear trace pairing
\((X,Y)\mapsto\Tr(XY)\) is nondegenerate on \(M_2(\C)\).  Every entry of
the Gram matrix for this basis belongs to \(k=\Q(\Tr\Delta)\), because
Cayley--Hamilton rewrites powers of \(A\) and \(C\) using their traces and
all remaining products are group elements whose traces lie in \(k\).
Pairing an arbitrary element of \(\Delta\) against the four basis elements
and solving the resulting linear system therefore gives coefficients in
\(k\).  Products of elements of \(\Delta\) are again in \(\Delta\), so its
\(k\)-span is closed under multiplication and has this four-element basis.
After extension to \(\C\) it is \(M_2(\C)\); it is consequently a
four-dimensional central simple \(k\)-algebra, or equivalently a
quaternion algebra over \(k\).
\end{proof}

The \(k\)-subalgebra of \(M_2(\R)\) generated by \(\Delta\) is a
quaternion algebra \(B/k\), because \(\Delta\) is nonelementary.  After
enlarging a finite set \(S\) of finite places, there is a smooth affine
group scheme \(\mathcal H/\cO_{k,S}\), with generic fiber
\(\SL_1(B)\), such that
\begin{equation}\label{eq:integral-model}
 \Delta\subset\mathcal H(\cO_{k,S}).
\end{equation}
Indeed, choose an \(\cO_k\)-order, or merely a full
\(\cO_k\)-lattice, in \(B\).  Enlarge \(S\) so that the finitely many
chosen generators of \(\Delta\) lie in its
\(\cO_{k,S}\)-localization, and then remove the finitely many further
places at which the resulting reduced-norm-one model is not smooth.
We always enlarge \(S\) by all finite places over a rational prime whenever
one place over that prime is included.  Such an \(S\) will be called
\emph{rationally saturated}.  We also include all places above the rational
primes ramified in \(k\).  Thus every rational prime outside \(S\) is
unramified in \(k\).  Norms of \(S\)-ideals are the absolute norms defined
in \eqref{eq:ideal-norm-definition}; equivalently, only their
prime-to-\(S\) factors contribute.

The restriction-of-scalars group associated with \eqref{eq:integral-model}
is
\[
 \mathbf G=\Res_{k/\Q}\SL_1(B).
\]
By definition, restriction of scalars is the \(\Q\)-group characterized by
\[
 \mathbf G(R)=\SL_1(B)(k\otimes_\Q R)
\]
for every \(\Q\)-algebra \(R\).  After extension to \(\overline\Q\), it
becomes one copy of \(\SL_2\) for every embedding
\(k\hookrightarrow\overline\Q\).

\begin{lemma}[Zariski closure after restriction of scalars]
\label{lem:full-restriction-closure}
The connected \(\Q\)-Zariski closure of the restriction-of-scalars image of
\(\Delta\) is \(\mathbf G\).
\end{lemma}

\begin{proof}
The group \(\Delta\) is nonelementary, hence Zariski dense in
\(\SL_1(B)\) over \(k\).  Over \(\overline\Q\), the group \(\mathbf G\)
is a product of \(\SL_2\)-factors indexed by the embeddings
\(\sigma:k\hookrightarrow\overline\Q\), and the connected closure of
\(\Delta\) projects onto every factor.

A connected subgroup that projects onto every factor of a product of
simple groups of type \(A_1\) is a product of diagonal subgroups.  Here is
the relevant consequence of Goursat's lemma.  For any two factors, the
image of the subgroup in their product is either the full product or the
graph of an isomorphism between the factors.  At the level of Lie algebras,
this follows because the kernel of either projection is an ideal in
\(\mathfrak{sl}_2\), and hence is either zero or all of
\(\mathfrak{sl}_2\).  Declare two factors equivalent when their two-factor
image is such a graph.  The subgroup is then the product, over the
equivalence classes, of copies of \(\SL_2\) embedded diagonally in the
factors belonging to that class.  Suppose that one such diagonal subgroup
identified the factors corresponding to two different embeddings
\(\sigma\ne\tau\).  On a finite-index subgroup \(\Delta_0<\Delta\), the
two representations would then differ by an algebraic automorphism of
\(\SL_2\).  Every such automorphism preserves trace, and consequently
\[
 \sigma(\Tr h)=\tau(\Tr h)\qquad(h\in\Delta_0).
\]
The invariant trace field of a nonelementary Fuchsian group is unchanged
on passage to finite index; see
\cite[Theorem~3.3.4]{MaclachlanReid}.  Hence
\[
 k=\Q\bigl(\Tr(h^2):h\in\Delta_0\bigr).
\]
The displayed equality of traces for \(h\in\Delta_0\), applied also to
\(h^2\), makes \(\sigma\) and \(\tau\) agree on this generating set of
\(k\), forcing \(\sigma=\tau\).  Thus every diagonal block is a singleton,
which proves the assertion.
\end{proof}

The choice of the full trace field is essential here.  Over a larger
coefficient field, the Zariski closure could identify distinct
restriction-of-scalars factors along a diagonal subgroup.  Its congruence
quotients would then have smaller dimension, invalidating the quotient-size
estimate used below.

\section{Local expansion and the Chebyshev recurrence}
\label{sec:local-expansion}

\subsection{The local expansion functional}

Let \(k\) be a number field, and write \(M_k\) for its set of archimedean
and nonarchimedean places.  Absolute values are normalized to extend the
usual absolute values on \(\Q\), and \(n_v=[k_v:\Q_v]\) denotes the local
degree.  For \(t\in k\) and a place \(v\), choose a root
\(z_v\) in an algebraic closure of \(k_v\) of
\begin{equation}\label{eq:z-equation}
 z+z^{-1}=t^2-2.
\end{equation}
Define
\begin{equation}\label{eq:local-chi}
 \chi_v(t)=\frac12\log\max\{|z_v|_v,|z_v|_v^{-1}\},
 \qquad
 \cJ_k(t)=\sum_{v\in M_k}n_v\chi_v(t).
\end{equation}
The two roots are reciprocal, so replacing \(z_v\) by \(z_v^{-1}\) does
not change \(\chi_v(t)\); the extension of the absolute value to the finite
splitting field is unique up to conjugacy and gives the same value.  The
element \(t\) is integral at all but finitely many finite places.  At every
remaining finite place both roots are units, and hence only finitely many
nonarchimedean terms in \(\cJ_k(t)\) are nonzero.

Our absolute logarithmic Weil height is normalized by
\begin{equation}\label{eq:height-normalization}
 [k:\Q]h(t)=\sum_{v\in M_k}n_v\log^+|t|_v,
 \qquad \log^+r=\max\{0,\log r\}.
\end{equation}
The product formula in the same normalization is
\(\sum_vn_v\log|a|_v=0\) for \(a\in k^\times\).

\begin{lemma}[basic properties and change of field]
\label{lem:J-base-change}
Every local summand in \eqref{eq:local-chi} is nonnegative, and
\(\cJ_k(-t)=\cJ_k(t)\).  If \(L/k\) is a finite extension and \(t\in k\),
then
\begin{equation}\label{eq:J-base-change}
 \cJ_L(t)=[L:k]\cJ_k(t).
\end{equation}
\end{lemma}

\begin{proof}
The first two assertions follow directly from the definition and from the
fact that \eqref{eq:z-equation} depends only on \(t^2\).  If \(w\) is a
place of \(L\) above \(v\), the normalized absolute value restricts to the
one at \(v\), so \(\chi_w(t)=\chi_v(t)\).  The local-degree formula is
\(\sum_{w\mid v}n_w=[L:k]n_v\).
Summing over \(v\) proves \eqref{eq:J-base-change}.
\end{proof}

The definition can also be expressed in terms of eigenvalues.  If
\(\lambda+\lambda^{-1}=t\),
then \(z=\lambda^2\) satisfies \eqref{eq:z-equation}, and hence
\begin{equation}\label{eq:lambda-chi}
 \chi_v(t)=\log\max\{|\lambda|_v,|\lambda|_v^{-1}\}.
\end{equation}
At a real Fuchsian place, if \(|t|=2\cosh(L/2)\), one may take
\(|\lambda|=e^{L/2}\), so
\begin{equation}\label{eq:distinguished-contribution}
 \chi_{v_0}(t)=\frac L2.
\end{equation}

\begin{lemma}[comparison with Weil height]
\label{lem:J-height-comparison}
For every number field \(k\) and every \(t\in k\),
\begin{equation}\label{eq:J-height-comparison}
 \left|\cJ_k(t)-[k:\Q]h(t)\right|
 \le [k:\Q]\log2.
\end{equation}
At a nonarchimedean place the two local contributions agree exactly.
\end{lemma}

\begin{proof}
Choose \(\lambda\) as in \eqref{eq:lambda-chi} and put
\(M_v=\max\{|\lambda|_v,|\lambda|_v^{-1}\}\ge1\).  At a
nonarchimedean place, the ultrametric inequality applied to
\(t=\lambda+\lambda^{-1}\) shows that \(|t|_v=M_v\) when \(M_v>1\),
whereas \(|t|_v\le1\) when \(M_v=1\).  Thus
\(\log M_v=\log^+|t|_v\).

At an archimedean place,
\[
 |t|_v\le |\lambda|_v+|\lambda|_v^{-1}\le2M_v,
\]
while the quadratic formula gives
\[
 M_v\le |t|_v+1\le2\max\{1,|t|_v\}.
\]
Therefore
\[
 \left|\log M_v-\log^+|t|_v\right|\le\log2.
\]
Multiply by \(n_v\), sum over all places, and use
\(\sum_{v\mid\infty}n_v=[k:\Q]\).
\end{proof}

In particular, if \(\cJ_k(t)=0\), then \(t\) is an algebraic integer and
every archimedean conjugate of \(t\) belongs to \([-2,2]\).  More
generally, if the entire sum \(\cJ_k(t)\) comes from one real place, then
every finite-place term vanishes, so \(t\) is an algebraic integer, and every
other archimedean conjugate belongs to \([-2,2]\).

\subsection{Recurrence identities}

Define \(q_n\in\Z[T]\) by
\begin{equation}\label{eq:q-recurrence}
 q_0=0,\qquad q_1=1,\qquad q_{n+1}=Tq_n-q_{n-1}.
\end{equation}
Thus \(q_n(T)=U_{n-1}(T/2)\), where \(U_j\) is the Chebyshev polynomial of
the second kind.

\begin{lemma}[matrix recurrence and Cassini identity]
\label{lem:matrix-recurrence}
Let \(A,G\in\SL_2(K)\), where \(K\) is a field, and put
\(s=\Tr A\).  For every \(n\ge1\),
\begin{align}
 A^n&=q_n(s)A-q_{n-1}(s)I,
 \label{eq:matrix-recurrence}\\
 \Tr(A^nG)&=q_n(s)\Tr(AG)-q_{n-1}(s)\Tr G.
 \label{eq:trace-recurrence}
\end{align}
Moreover,
\begin{equation}\label{eq:cassini}
 q_n(T)^2-q_{n-1}(T)q_{n+1}(T)=1.
\end{equation}
Consequently two consecutive specialized values
\(q_{n-1}(s),q_n(s)\) generate the unit ideal in every ring containing
\(s\).
\end{lemma}

\begin{proof}
Cayley--Hamilton gives \(A^2-sA+I=0\).  Induction using
\eqref{eq:q-recurrence} proves \eqref{eq:matrix-recurrence}; multiplication
by \(G\) and taking traces gives \eqref{eq:trace-recurrence}.  The left side
of \eqref{eq:cassini} is independent of \(n\) by the recurrence and equals
one when \(n=1\).  The final assertion follows directly from
\eqref{eq:cassini}.
\end{proof}

Two forms of \eqref{eq:trace-recurrence} will be used repeatedly.  Put
\(P=q_{N-1}(s)\), \(Q=q_N(s)\).  Then
\begin{align}
 Q\Tr(AG)-P\Tr G&=\Tr(A^NG),
 \label{eq:positive-recurrence-image}\\
 Q\Tr G-P\Tr(AG)&=\Tr(A^{-(N-1)}G).
 \label{eq:negative-recurrence-image}
\end{align}
The second identity follows by applying Cayley--Hamilton backwards, or by
multiplying \eqref{eq:matrix-recurrence} by \(A^{-N}\).

Because \(q_j\in\Z[T]\), every \(q_j(t)\) is \(S\)-integral whenever
\(t\in\cO_{k,S}\).  In the principal argument the recurrence index \(N\)
is fixed and the trace varies; no asymptotic assertion in the recurrence
index is needed.

\section{External results used in the proof}
\label{sec:external-inputs}

This section records the analytic and arithmetic theorems used below, with
their required uniformity made explicit.  Standard geometric or algebraic
facts used only once are stated where they enter the proof.  The deductions
from these results are proved in the subsequent sections.

\subsection{Diophantine approximation and height finiteness}

We begin by recalling the local notation in Roth's theorem.  If
\(P=(x_0:x_1)\) and \(Q=(y_0:y_1)\) are points of
\(\mathbb P^1(\overline{k_v})\), their chordal distance is
\[
 \delta_v(P,Q)=
 \frac{|x_0y_1-x_1y_0|_v}
      {\lVert(x_0,x_1)\rVert_v\lVert(y_0,y_1)\rVert_v},
\]
where the maximum norm is used at a nonarchimedean place and a Euclidean
norm at an archimedean place.  A local projective proximity function is
\(\lambda_{Q,v}(P)=-\log\delta_v(P,Q)\).  Changing these norms changes the
function by a bounded amount, which is absorbed by the \(O(1)\) term below.
When the target is algebraic over \(k\), one first fixes an extension of
\(|\cdot|_v\) to \(\overline k\), or equivalently a compatible embedding
\(\iota_v:\overline k\hookrightarrow\overline{k_v}\), and applies the
same formula to \(\iota_v(Q)\).  This choice is part of the local target
data; no canonical embedding is being asserted.

\begin{theorem}[simultaneous-place Roth theorem]
\label{input:roth}
Let \(k\) be a number field, let \(T\) be a finite set of places, and at
each \(v\in T\) choose one algebraic target
\(\alpha_v\in\mathbb P^1(\overline k)\), with \(\alpha_v=\infty\) allowed,
and a compatible embedding
\(\iota_v:\overline k\hookrightarrow\overline{k_v}\).
For every \(\epsilon>0\), outside a finite set of \(t\in k\),
\[
 \sum_{v\in T}n_v\lambda_{\iota_v(\alpha_v),v}(t)
 \le(2+\epsilon)[k:\Q]h(t)+O(1),
\]
where \(t\) is regarded as the projective point \((t:1)\) and
\(\lambda_{\iota_v(\alpha_v),v}\) is the local proximity function just
defined.
\end{theorem}

Theorem~\ref{input:roth} is the projective form of
\cite[Theorem~6.4.1]{BombieriGubler}, whose targets are written in an
affine coordinate.  To pass to the displayed form, choose
\(\beta\in k\) different from every finite target and put
\(\phi(z)=(z-\beta)^{-1}\).  Then every \(\phi(\alpha_v)\) is finite,
including \(\phi(\infty)=0\).  Apart from the single point
\(t=\beta\), which may be added to the exceptional set, functoriality of
local Weil functions and heights under this fixed projective transformation
gives
\[
 \lambda_{\alpha_v,v}(t)
 =\lambda_{\phi(\alpha_v),v}(\phi(t))+O_v(1),
 \qquad h(\phi(t))=h(t)+O_\phi(1).
\]
Here the chosen embeddings into the local algebraic closures are retained
when \(\phi\) is applied.  Applying the cited affine theorem to \(\phi(t)\)
and summing the finitely many bounded local errors proves the stated
projective version, including a target at infinity.
Section~\ref{sec:roth-sieve} applies this theorem to a rational function of
degree \(N-1\), producing the coefficient \(2N-4\).

\begin{theorem}[Northcott finiteness theorem]
\label{input:northcott}
For fixed \(D\) and \(H\), the set of algebraic numbers \(\alpha\) with
\([\Q(\alpha):\Q]\le D\) and \(h(\alpha)\le H\) is finite.
\end{theorem}

See \cite[Theorem~1.6.8]{BombieriGubler}.

\subsection{Congruence expansion and counting}

We shall use congruence expansion in two forms.  The random-walk estimate is
used to choose a short perturbation of a prescribed group element.  The
geometric-ball estimate is used later to count trace classes.  Both follow
from strong approximation, superapproximation for bounded prime powers,
the transfer of a spectral gap from Cayley graphs to surface covers, and
uniform lattice-point counting.

Let \(y\) be an algebraic Fuchsian specialization, let
\(\Delta=\rho_y(\Pi^{(2)})\), and retain the smooth
\(\cO_{k_y,S}\)-model \(\mathcal H\) of \(\SL_1(B_y)\) from
\eqref{eq:integral-model}.  The restriction of scalars
\(\mathbf G=\Res_{k_y/\Q}\SL_1(B_y)\) is an algebraic group over \(\Q\).
By contrast, the finite groups used for congruence counting are
\(\mathcal H(\cO_{k_y,S}/\mathfrak c)\).
Lemma~\ref{lem:full-restriction-closure} verifies that the
connected \(\Q\)-Zariski closure is the connected, semisimple, simply
connected group \(\mathbf G\).  Congruence reduction is applied to the
linear group \(\Delta<\SL_2\).  Its congruence kernels, transported through
\(\Delta\simeq\overline\Delta\), are finite-index subgroups of the
projective surface lattice and therefore define covers of the fixed surface.
The optional central element in the full inverse image affects indices by at
most two and is absorbed in fixed constants.

Let \(q_0\) be the product of the rational primes below \(S\).  Rational
saturation gives \(\cO_{k_y,S}=\cO_{k_y}[1/q_0]\).  Form the smooth affine
\(\Z[1/q_0]\)-group scheme
\[
 \mathcal G_{\mathrm{res}}=
 \Res_{\cO_{k_y,S}/\Z[1/q_0]}\mathcal H.
\]
Its generic fiber is \(\mathbf G\), and for every integer \(m\) prime to
\(q_0\),
\[
 \mathcal G_{\mathrm{res}}(\Z[1/q_0]/m\Z[1/q_0])
 =\mathcal H(\cO_{k_y,S}/m\cO_{k_y,S}).
\]
Choose a faithful \(\Q\)-representation of \(\mathbf G\).  After enlarging
\(S\), spreading out gives a closed immersion
\[
 \iota:\mathcal G_{\mathrm{res}}\hookrightarrow\GL_{M,\Z[1/q_0]}
\]
for one fixed integer \(M\).  Consequently
\(\iota(\Delta)\subset\GL_M(\Z[1/q_0])\), and reduction of this matrix
group modulo \(m\) is exactly the reduction of \(\Delta\) modulo
\(m\cO_{k_y,S}\).  We use this integral model to apply strong
approximation and superapproximation to the same congruence quotients.

\begin{theorem}[strong approximation in the form used here]
\label{input:strong-approximation}
After deleting finitely many rational primes, there is a constant
\(C_{\mathrm{sa}}\) such that the image \(Q_{\mathfrak c}\) of
\(\Delta\) modulo every prime-to-\(S\) ideal \(\mathfrak c\) satisfies
\begin{equation}\label{eq:strong-approximation-index}
 [\mathcal H(\cO_{k_y,S}/\mathfrak c):Q_{\mathfrak c}]
 \le C_{\mathrm{sa}}.
\end{equation}
Moreover,
\begin{equation}\label{eq:quotient-size}
 \#Q_{\mathfrak c}\asymp_{y,S}(\Norm\mathfrak c)^3,
\end{equation}
uniformly in \(\mathfrak c\).
\end{theorem}

We explain why strong approximation applies over the coefficient ring used
in the congruence quotients.  Put
\[
 A_{\mathrm{ad}}
 =\Z\bigl[\Tr(\operatorname{Ad}h):h\in\Delta\bigr]\subset k_y.
\]
For \(h\in\SL_2\),
\[
 \Tr(\operatorname{Ad}h)=(\Tr h)^2-1=\Tr(h^2)+1.
\]
The invariant-trace-field identity, together with its invariance under
finite-index passage, therefore gives
\(\operatorname{Frac}(A_{\mathrm{ad}})=k_y\).  Because \(\Delta\) is
finitely generated, finite generation of matrix trace algebras
\cite{Procesi,Horowitz} makes \(A_{\mathrm{ad}}\) a finitely generated
trace ring.
After inverting finitely many rational primes, it is the localization of an
order in \(k_y\).  Enlarge the rationally saturated set \(S\) to contain
every prime ideal dividing the conductor of that order.  Localization then
identifies the order with the full ring of \(S\)-integers, so
\[
 A_{\mathrm{ad}}[1/q_0]=\cO_{k_y,S}.
\]

Apply \cite[Theorem~9.1.1]{WeisfeilerStrong}; see also
\cite{MVW,Nori}.  That theorem first supplies a localization of
\(A_{\mathrm{ad}}\) and a smooth simply connected semisimple model
\(\mathcal H_{\mathrm W}\) of the \(k_y\)-group \(\SL_1(B_y)\), for which
the closure of \(\Delta\) in the product of the retained local integral
groups is open.  Enlarge \(S\) to contain the primes occurring in that
localization.  The preceding equality then permits base change of
\(\mathcal H_{\mathrm W}\) to \(\cO_{k_y,S}\).  Its Weil restriction to
\(\Z[1/q_0]\) and \(\mathcal G_{\mathrm{res}}\) both have generic fiber
\[
 \mathbf G=\Res_{k_y/\Q}\SL_1(B_y).
\]
After one further localization, a generic-fiber isomorphism and its inverse
extend to an isomorphism between these two models.  Thus Weisfeiler's model,
initially defined from the trace ring, agrees with the congruence model
\(\mathcal H\) outside the enlarged finite set \(S\).
Lemma~\ref{lem:full-restriction-closure} verifies that the
connected \(\Q\)-Zariski closure is all of the simply connected semisimple
group \(\mathbf G\).  Consequently the closure
\(\widehat\Delta\) of \(\Delta\) is an open subgroup of
\[
 K=\prod_{p\nmid q_0}\mathcal G_{\mathrm{res}}(\Z_p)
  =\prod_{p\nmid q_0}
    \mathcal H(\cO_{k_y}\otimes_\Z\Z_p).
\]
Put \(C_{\mathrm{sa}}=[K:\widehat\Delta]\).  For every prime-to-\(S\)
ideal \(\mathfrak c\), smoothness and the Chinese remainder theorem make
the reduction map
\[
 r_{\mathfrak c}:K\longrightarrow
 \mathcal H(\cO_{k_y,S}/\mathfrak c)
\]
surjective.  Since the target is finite and \(\Delta\) is dense in
\(\widehat\Delta\),
\(r_{\mathfrak c}(\widehat\Delta)=r_{\mathfrak c}(\Delta)=Q_{\mathfrak c}\).
For a surjective homomorphism from a compact group to a finite group, the
index of the image of an open subgroup is at most the index of that open
subgroup.  Hence
\[
 [\mathcal H(\cO_{k_y,S}/\mathfrak c):Q_{\mathfrak c}]
 \le [K:\widehat\Delta]=C_{\mathrm{sa}},
\]
which proves \eqref{eq:strong-approximation-index} simultaneously for every
prime-to-\(S\) ideal, with no restriction on its prime exponents.

It remains to calculate the size of the finite quotient.  At a good prime
\(\mathfrak p\), with \(q=\Norm\mathfrak p\), smooth lifting in relative
dimension three and the fact that every quaternion algebra over a finite
field splits give
\[
 \#\mathcal H(\cO_{k_y}/\mathfrak p^e)
 =q^{3(e-1)}\#\SL_2(\mathbb F_q)
 =q^{3e}(1-q^{-2}).
\]
The Chinese remainder theorem and
\(\prod_{\mathfrak p}(1-(\Norm\mathfrak p)^{-2})>0\) show that
\(\#\mathcal H(\cO_{k_y,S}/\mathfrak c)
\asymp(\Norm\mathfrak c)^3\).  Equation
\eqref{eq:strong-approximation-index} now proves
\eqref{eq:quotient-size}.  This use of strong approximation provides only
the bounded index.  The spectral gap required later comes from
superapproximation.

Fix a finite symmetric generating set \(\Omega\) of \(\Delta\).  The
\emph{lazy random walk} used below starts at the identity and, independently
at each step, stays in place with probability \(1/2\) and otherwise
right-multiplies by a uniformly chosen element of \(\Omega\).  We denote its
position after \(L\) steps by \(W_L\).  Thus a walk with starting point
\(g_0\) ends at \(g_0W_L\).  Laziness makes the transition operator
aperiodic and removes a possible spectral obstruction at the eigenvalue
\(-1\).

Under the closed immersion \(\iota\) defined above,
\(\Omega\subset\GL_M(\Z[1/q_0])\), and
Lemma~\ref{lem:full-restriction-closure} says that the identity component
of its \(\Q\)-Zariski closure is \(\iota(\mathbf G)\), hence is perfect.
Thus \cite[Theorem~1]{SalehiSuperII} applies to the reductions of this
matrix group modulo \(M_0^D\), with \(M_0\) square-free and prime to
\(q_0\).  By the defining property of \(\mathcal G_{\mathrm{res}}\), the
resulting reduction image is
\[
 Q_{M_0^D\cO_{k_y,S}}.
\]
If \(\mathfrak c\) has all prime-ideal exponents at most \(D\), let
\(M_0\) be the product of the rational primes below \(\mathfrak c\).  All
such primes are unramified because the ramified rational primes were put in
\(S\); consequently
\[
 M_0^D\cO_{k_y,S}\subseteq\mathfrak c.
\]
Reduction therefore induces a surjection
\(Q_{M_0^D\cO_{k_y,S}}\twoheadrightarrow Q_{\mathfrak c}\).  For the
symmetric lazy walk used here, functions on \(Q_{\mathfrak c}\) pull back
along the quotient map to an invariant subspace of functions on
\(Q_{M_0^D\cO_{k_y,S}}\).  Hence every nonconstant eigenvalue of the
quotient walk already occurs upstairs, so the same normalized spectral-gap
lower bound holds for \(Q_{\mathfrak c}\).

\begin{theorem}[mixing modulo prime squares]
\label{input:prime-square-mixing}
Use the generating set and lazy random walk defined above.  After removing finitely
many rational primes, there are constants \(B,C,\delta>0\) such that, for
every remaining rational prime \(p\), every starting point \(g_0\), and
every subset \(\mathcal A\) of the reduction image modulo \(p^2\),
\begin{equation}\label{eq:prime-square-mixing}
 \Pr(g_0W_L\bmod p^2\in\mathcal A)
 \le C\frac{\#\mathcal A}{\#Q_{p^2}}+p^B e^{-\delta L}.
\end{equation}
Here reduction modulo \(p^2\) means reduction modulo
\(p^2\cO_{k_y,S}\), and
\(Q_{p^2}=Q_{p^2\cO_{k_y,S}}\) is the image of \(\Delta\) in that
quotient.  The constants are independent of \(p,L,g_0\).
\end{theorem}

The bounded-power superapproximation theorem just discussed, applied with
exponent \(2\), supplies a uniform spectral gap for these congruence
quotients.  Making the walk lazy removes a
possible eigenvalue near \(-1\).  Starting from a point mass, the resulting
\(L^2\)-distance from the uniform distribution, with \(L^2\) formed using
unnormalized counting measure on the finite quotient, is
\(O(e^{-\delta L})\).  Cauchy--Schwarz bounds its mass on \(\mathcal A\) by
\[
 \sqrt{\#\mathcal A}\,O(e^{-\delta L})
 \le\sqrt{\#Q_{p^2}}\,O(e^{-\delta L})
 \le p^B e^{-\delta L};
\]
the last inequality follows from the dimension-three local count above.
This proves \eqref{eq:prime-square-mixing}.  Only the fixed exponent two is
used here.

\begin{theorem}[congruence counting for moduli with bounded prime
exponents]
\label{input:bounded-exponent-counting}
Fix a faithful cocompact Fuchsian character \(x\), an algebraic faithful
Fuchsian character \(y\), and an integer \(D\ge1\).  Let
\[
 \cB_R^x=\{g\in\Pi^{(2)}:d_x(o,go)\le R\}.
\]
There is a constant \(\delta_D>0\) with the following property.  If
\(\mathfrak c\subset\cO_{k_y,S}\) is prime to \(S\) and every prime
exponent of \(\mathfrak c\) is at most \(D\), then, in the range
\(\log\Norm\mathfrak c=O(R)\) and uniformly in a trace class
\(u\bmod\mathfrak c\),
\begin{equation}\label{eq:bounded-exponent-trace-count}
 \#\{g\in\cB_R^x:\Tr_y g\equiv u\pmod{\mathfrak c}\}
 \le
 \frac{e^{R+o(R)}}{\Norm\mathfrak c}
 +(\Norm\mathfrak c)^{2+o(1)}e^{(1-\delta_D)R}.
\end{equation}
The error terms are uniform for fixed \(x,y,D\).  In particular, the second
term is absorbed by the first whenever
\begin{equation}\label{eq:counting-absorption-condition}
 3\log\Norm\mathfrak c<\delta_D R-o(R).
\end{equation}
Consequently, if \(\log\Norm\mathfrak c_m\le Am+O(1)\), choose, after
\(D,A,\delta_D\) have been fixed, a constant
\begin{equation}\label{eq:radius-choice}
 C>\frac{3A}{\delta_D}
\end{equation}
and put \(R=Cm\).  Then
\begin{equation}\label{eq:bounded-exponent-flat-form}
 \max_u\#\{g\in\cB_R^x:\Tr_y g\equiv u\pmod{\mathfrak c_m}\}
 \le\frac{e^{R+o(m)}}{\Norm\mathfrak c_m}.
\end{equation}
\end{theorem}

In this statement a \emph{trace class} is an element
\(u\in\cO_{k_y,S}/\mathfrak c\), and
\(\Tr_y g\equiv u\pmod{\mathfrak c}\) means that the trace of the integral
matrix representing \(g\) maps to \(u\) under the quotient homomorphism.
The displayed \(o(R)\)-terms are uniform in \(\mathfrak c\) and \(u\) once
\(x,y,D\) and a bound for \(\log\Norm\mathfrak c/R\) have been fixed.

The integral-model argument preceding
Theorem~\ref{input:prime-square-mixing} gives a uniform spectral gap for
the Cayley graphs of all the quotients occurring in this theorem.  We next
explain how this discrete spectral gap yields the stated count in a
hyperbolic displacement ball.  Let
\(\bar\rho_x:\Pi^{(2)}\to\PSL_2(\R)\) be the projective Fuchsian
representation with character \(x\).  Since
\(\rho_y:\Pi^{(2)}\to\Delta\) is an isomorphism onto its image, put
\[
\begin{aligned}
 G_\infty&=\PSL_2(\R),
 &\Lambda_x&=\bar\rho_x(\Pi^{(2)}),\\
 \pi_{\mathfrak c}&:\Delta\longrightarrow Q_{\mathfrak c},
 &K_{\mathfrak c}&=\ker(\pi_{\mathfrak c}),\\
 \bar\pi_{\mathfrak c}
 &=\pi_{\mathfrak c}\circ\rho_y:
      \Pi^{(2)}\longrightarrow Q_{\mathfrak c},\\
 \Pi_{\mathfrak c}&=\ker(\bar\pi_{\mathfrak c}),
 &\Lambda_{x,\mathfrak c}&=\bar\rho_x(\Pi_{\mathfrak c}).
\end{aligned}
\]
Let \(m_\infty\) be Haar measure on \(G_\infty\), and let
\(B_R\subset G_\infty\) be the set of elements moving \(o\) by at most
\(R\).
Because \(K_{\mathfrak c}\) is the kernel of a quotient map, its pullback
\(\Pi_{\mathfrak c}\) is normal in \(\Pi^{(2)}\).  Thus the
associated surface cover is regular (Galois).  Its deck group is
\(Q_{\mathfrak c}\), and its quotient graph for the fixed generating set
is the Cayley graph whose spectral gap was obtained above.
The Brooks--Burger transfer \cite{BrooksTower,Burger}, specifically
Burger's \cite[Th\'eor\`eme~B]{Burger}, gives a number \(\lambda_D>0\),
independent of \(\mathfrak c\), such that
\[
 \lambda_1(\Lambda_{x,\mathfrak c}\backslash\mathbb H^2)\ge\lambda_D.
\]

Let \(p:\widetilde G=\SL_2(\R)\to G_\infty\) be the central covering, and
put
\[
 \widetilde\Lambda_x=p^{-1}(\Lambda_x),\qquad
 \widetilde\Lambda_{x,\mathfrak c}
     =p^{-1}(\Lambda_{x,\mathfrak c}).
\]
The natural maps
\[
 \widetilde G/\widetilde\Lambda_x\simeq G_\infty/\Lambda_x,
 \qquad
 \widetilde G/\widetilde\Lambda_{x,\mathfrak c}
       \simeq G_\infty/\Lambda_{x,\mathfrak c}
\]
identify the corresponding unitary representations and the averages over a
displacement ball with the averages over its full inverse image.  Inversion,
\(g\widetilde\Lambda\mapsto\widetilde\Lambda g^{-1}\), identifies the
right homogeneous spaces used here with the left-quotient convention for
surface covers and in the cited lattice-counting results.  Set
\[
 \lambda_D'=\min\{\lambda_D,1/4\},\qquad
 c_D=\frac{1-\sqrt{1-4\lambda_D'}}2>0.
\]
Let \(K_\infty=\operatorname{PSO}(2)\) be a maximal compact subgroup of
\(G_\infty\), and put
\(\widetilde K_\infty=p^{-1}(K_\infty)=\operatorname{SO}(2)\).
A nontempered irreducible constituent of
\(L^2_0(G_\infty/\Lambda_{x,\mathfrak c})\) is a complementary-series
representation with parameter \(s\in(1/2,1)\), Laplace eigenvalue
\(s(1-s)\), and \(K_\infty\)-finite matrix coefficients bounded, up to a
polynomial factor, by \(e^{-(1-s)t}\) on the Cartan ray.  The displayed
lower bound gives \(1-s\ge c_D\); see \cite{LangSL2} for the
complementary-series classification and \cite{CowlingHaagerupHowe} for the
almost-\(L^2\) matrix-coefficient criterion.  Since Haar measure in Cartan coordinates
has density \(\asymp e^t\,dt\), choosing
\[
 p_D>\max\{2,c_D^{-1}\}
\]
shows uniformly in \(\mathfrak c\) that the representations on
\(L^2_0(\widetilde G/\widetilde\Lambda_{x,\mathfrak c})\) are
\(L^{p_D+}\).  Here this notation means that every
\(\widetilde K_\infty\)-finite matrix coefficient belongs to
\(L^{p_D+\varepsilon}(\widetilde G)\) for every \(\varepsilon>0\).

The algebraic group \(\SL_2\) is simply connected and \(\R\)-isotropic,
so its group of real points is within the scope of
\cite[Theorem~4.5]{GorodnikNevoCounting}.  For the normalized averages
\(\widetilde\beta_R\) on the inverse images of the displacement balls, that
theorem gives, after choosing its auxiliary exponent sufficiently small,
\[
 \bigl\|\pi^0_{\mathfrak c}(\widetilde\beta_R)\bigr\|
 \le C_D\,m_\infty(B_R)^{-\kappa_D}
\]
where \(\pi^0_{\mathfrak c}\) denotes the averaging representation on the
mean-zero subspace of
\(L^2(\widetilde G/\widetilde\Lambda_{x,\mathfrak c})\).
Here \(C_D,\kappa_D>0\) are independent of \(\mathfrak c\); the same estimate
holds for the standard inner and outer thickenings of the balls.
Transporting this operator-norm bound through the two quotient
identifications above gives the same quantitative mean-ergodic estimate on
\(G_\infty\), with constants independent of \(\mathfrak c\).

Hyperbolic displacement balls are H\"older well rounded: thickening or
shrinking the boundary by a small fixed neighborhood changes their Haar
measure by a controlled power of the neighborhood size.  Thus
\cite[Theorem~1.9 and Corollaries~1.11, 1.13]
{GorodnikNevoCounting} turn the preceding bound into a lattice-point
estimate whose absolute error is \(O_D(e^{(1-\delta_D)R})\) for some
\(\delta_D>0\).  The injectivity neighborhood may be chosen uniformly:
each surface in the family covers the same fixed compact surface, so no
new shorter closed geodesic is introduced.  Moreover, a quotient
element has a representative in the fixed base lattice \(\Lambda_x\), and
its fiber is the corresponding translate of
\(\Lambda_{x,\mathfrak c}\).  The estimate is uniform over these translated
cosets by \cite[Remark~1.12]{GorodnikNevoCounting}.  Hence both the
congruence subgroup and the chosen element of its finite quotient may vary
in the count below.

With this normalization, uniform lattice counting gives, for every quotient
element \(q\),
\[
 \#\{g\in\cB_R^x:\bar\pi_{\mathfrak c}(g)=q\}
 =\frac{m_\infty(B_R)}
        {m_\infty(G_\infty/\Lambda_x)\#Q_{\mathfrak c}}
  +O_D(e^{(1-\delta_D)R});
\]
see \cite[Theorem~1.9 and Corollaries~1.11, 1.13]
{GorodnikNevoCounting}.  Here \(m_\infty(B_R)=e^{R+O_x(1)}\).  We also need
the size of a trace fiber in a finite quotient.  The following direct
matrix count in
\(A=\cO_{k_y}/\mathfrak p^e\) gives
\[
 \#\{g\in\SL_2(A):\Tr g=t\}\le(e+1)(\Norm\mathfrak p)^{2e}.
\]
To obtain it, write
\(g=\left(\begin{smallmatrix}a&b\\c&t-a\end{smallmatrix}\right)\).
The determinant condition is \(bc=a(t-a)-1\).  For each of the
\((\Norm\mathfrak p)^e\) choices of \(a\), partition \(b\) by its valuation
between \(0\) and \(e\); summing the corresponding numbers of solutions for
\(c\) gives at most \((e+1)(\Norm\mathfrak p)^e\) pairs \((b,c)\).
Write \(\mathfrak c=\prod_{\mathfrak p}\mathfrak p^{e_{\mathfrak p}}\).
At every retained good prime the chosen integral splitting identifies
\(\mathcal H(\cO_{k_y,\mathfrak p}/\mathfrak p^e)\) with
\(\SL_2(\cO_{k_y,\mathfrak p}/\mathfrak p^e)\), and reduced trace becomes
ordinary matrix trace.  Thus the preceding calculation and the Chinese
remainder theorem give
\[
 \#\{h\in\mathcal H(\cO_{k_y,S}/\mathfrak c):\Tr h=u\}
 \le (\Norm\mathfrak c)^2
       \prod_{\mathfrak p\mid\mathfrak c}(e_{\mathfrak p}+1).
\]
If \(e_{\mathfrak p}\le D\), then
\[
 \prod_{\mathfrak p\mid\mathfrak c}(e_{\mathfrak p}+1)
 \le (D+1)^{\omega(\mathfrak c)}
 = (\Norm\mathfrak c)^{o(1)}.
\]
Here \(\omega(\mathfrak c)\) is the number of distinct prime-ideal divisors
of \(\mathfrak c\).
The last estimate is uniform for fixed \(k_y,D\): for every \(\eta>0\),
the prime ideals with
\(\Norm\mathfrak p\le(D+1)^{1/\eta}\) contribute a fixed constant, while
for every remaining prime ideal
\(D+1\le(\Norm\mathfrak p)^\eta\).  Hence the product is at most
\(C_{k_y,D,\eta}(\Norm\mathfrak c)^\eta\).
Intersecting a trace fiber with \(Q_{\mathfrak c}\) can only decrease its
size, whereas Theorem~\ref{input:strong-approximation} gives
\(\#Q_{\mathfrak c}\asymp(\Norm\mathfrak c)^3\).  Thus a trace fiber in
\(Q_{\mathfrak c}\) has size at most
\((\Norm\mathfrak c)^{2+o(1)}\) and relative density at most
\((\Norm\mathfrak c)^{-1+o(1)}\).  Summing the quotient-element count over
that fiber proves \eqref{eq:bounded-exponent-trace-count}.  The quotient of
the error term by the main term is
\[
 (\Norm\mathfrak c)^{3+o(1)}e^{-\delta_D R},
\]
which tends to zero in the range
\(3\log\Norm\mathfrak c<\delta_D R-o(R)\).

The constants are not claimed to be uniform as \(D\to\infty\).  In each
application, \(D\) is fixed first and the radius constant is then chosen to
satisfy the strict inequality \eqref{eq:radius-choice}.

\subsection{Algebraic geometry over finite fields}

\begin{theorem}[Lang--Weil consequence]
\label{input:lang-weil}
Let \(\mathcal G\) be a smooth affine group scheme over
\(\Z[1/M]\) whose geometric fibers are connected, and let
\(0\ne\Phi\in\Q[\mathcal G_\Q]\).  After enlarging \(M\), so that every
remaining fiber is geometrically integral, \(\Phi\) extends to the model,
and the zero locus of its reduction has dimension at most
\(\dim\mathcal G_{\mathbb F_p}-1\), one has
\[
 \frac{\#\{g\in\mathcal G(\F_p):\Phi(g)=0\}}
      {\#\mathcal G(\F_p)}
 =O_\Phi(p^{-1})
\]
for every \(p\nmid M\).  The same bound, up to a fixed constant, holds in a
subgroup of uniformly bounded index.
\end{theorem}

Indeed, if \(d=\dim\mathcal G_{\mathbb F_p}\), Lang--Weil gives
\(\#\mathcal G(\mathbb F_p)=p^d+O(p^{d-1/2})\), uniformly after the fixed
model has been chosen.  Applied to the finitely many top-dimensional
components of the zero locus, whose dimensions are at most \(d-1\), it
gives \(O_\Phi(p^{d-1})\) points there.  Division proves the displayed
ratio.  If \(H_p<\mathcal G(\mathbb F_p)\) has index at most \(C\), then
\[
 \frac{\#(H_p\cap V(\Phi))}{\#H_p}
 \le C\,\frac{\#V(\Phi)(\mathbb F_p)}
              {\#\mathcal G(\mathbb F_p)},
\]
which proves the bounded-index assertion.  This is the only use of
Lang--Weil: it permits one fixed prime of good reduction to be chosen so
that a specified nonzero regular function is nonvanishing on more than
three quarters of the congruence image; see \cite{LangWeil}.

\subsection{Hyperbolic counting and arithmeticity}

We use the following geometric facts for a fixed cocompact Fuchsian metric.

\begin{description}[style=nextline,leftmargin=2.7cm]
\item[Orbit counting]
For a finite-index subgroup of a closed surface group,
\(\#\cB_R^x=e^{R+O_x(1)}\).
\item[Metric comparison]
Word length and displacement in the fixed metric are quasi-isometric.
\item[Centralizers]
The centralizer of a nonidentity element is infinite cyclic, and a cocompact
torsion-free Fuchsian group has no nontrivial parabolic element.
\end{description}

To verify the orbit estimate, choose a compact fundamental domain of
diameter \(D_0\).
The translates whose orbit points lie in a ball of radius \(R-D_0\) are
contained in the radius-\(R\) hyperbolic ball, while every translate meeting
that ball has orbit point in the radius-\(R+D_0\) ball.  Division by the
area of the fundamental domain and
\(\operatorname{area}(B_R)=2\pi(\cosh R-1)\) gives
\(\#\cB_R^x=e^{R+O_x(1)}\); the same proof applies to a finite-index
subgroup.  The word/displacement comparison is the
\v{S}varc--Milnor lemma for the proper cocompact action on
\(\mathbb H^2\).  Finally, every nonidentity element is hyperbolic:
elliptics are excluded by torsion-freeness and parabolics by compactness of
the quotient.  A commuting element preserves its unique axis, and
discreteness makes the group of translations along that axis cyclic.  This
proves the centralizer assertion.

\begin{theorem}[finite-covolume corollary of Takeuchi's criterion]
\label{input:takeuchi}
Let \(\Lambda<\PSL_2(\R)\) be a finite-covolume Fuchsian group, let
\(\widetilde\Lambda\) be its full inverse image in \(\SL_2(\R)\), and let
\(k=\Q(\Tr\widetilde\Lambda)\) be its ordinary trace field.  Put
\[
 K_2=\Q\bigl((\Tr g)^2:g\in\widetilde\Lambda\bigr),
\]
the invariant trace field.  Suppose:
\begin{enumerate}[label=\textup{(\roman*)}]
\item \(k\) is a number field;
\item every trace in \(\widetilde\Lambda\) is an algebraic integer;
\item for every embedding \(\sigma:k\hookrightarrow\C\) different from
the distinguished real embedding and every
\(g\in\widetilde\Lambda\), one has
\(\sigma(\Tr g)\in[-2,2]\).
\end{enumerate}
Then \(\Lambda\) is derived from a quaternion algebra and is arithmetic.
\end{theorem}

This formulation is a corollary of Takeuchi's statements.  Takeuchi denotes
the ordinary trace field by \(k_1\) and the invariant trace field by
\(k_2\).  His Theorem~2, the criterion for being derived from a quaternion
algebra used here, imposes the conjugate-trace condition for every
nonidentity embedding of \(k_1\).  These are the embeddings covered by
hypothesis~(iii).  The condition involving embeddings of
\(k_1\) whose restriction to \(k_2\) is nonidentity occurs instead in
Takeuchi's Theorem~1, his arithmetic criterion.  Proposition~2 identifies
the cocompact case with total reality and the interval \([-2,2]\)
\cite[Theorem~2 and Proposition~2]{Takeuchi}.  Theorem~2 now gives
containment in the projective norm-one group associated with an order in
the resulting admissible quaternion algebra.  Both that norm-one group and
\(\Lambda\) are lattices.
Comparing covolumes shows that the containment has finite index; this is
the asserted arithmeticity.  The finite-covolume hypothesis is essential
for this last step: an infinite-index subgroup can have the full boundary
circle as its limit set without being a lattice.  We use the corollary with
\(\Lambda=\Gamma^{(2)}\), after proving explicitly that its ordinary trace
field is the invariant trace field of \(\Gamma\).  See also
\cite[Theorem~3.3.4 and Lemma~3.5.6]{MaclachlanReid}.

\section{A lower bound for the recurrence divisor}
\label{sec:roth-sieve}

We first use Roth's approximation theorem to bound from below the norm of
the product of two consecutive recurrence terms.

\begin{lemma}[reduced numerator--denominator divisor]
\label{lem:roth-reduced-divisor}
Let \(k\) be a number field, let \(f\in k(T)\) have degree \(d\ge2\), and
suppose every zero and pole of \(f\) on
\(\mathbb P^1_{\overline k}\) is simple.  Let \(S\) be a finite set of
finite places containing the places at which a fixed integral model of
\(f\), or the reduced divisor of its zeros and poles, has bad reduction.
For \(t\in k\) which is neither a zero nor a pole, put
\[
 \mathfrak m_{f,S}(t)=
 \prod_{\mathfrak p\notin S}
 \mathfrak p^{|\operatorname{ord}_{\mathfrak p}f(t)|}.
\]
Here \(\operatorname{ord}_{\mathfrak p}\) is the discrete valuation
normalized by \(\operatorname{ord}_{\mathfrak p}(\mathfrak p)=1\).  The
absolute value in the exponent records both numerator and denominator
primes, each with its full valuation.
For every \(\epsilon>0\), outside a finite set of \(t\in k\),
\begin{equation}\label{eq:roth-reduced-divisor}
 \log\Norm\mathfrak m_{f,S}(t)
 \ge(2d-2-\epsilon)[k:\Q]h(t)-O_{f,S,\epsilon}(1).
\end{equation}
\end{lemma}

\begin{proof}
Let \(T_0\) consist of \(S\) and all archimedean places.  The product
formula and functoriality of height give
\begin{equation}\label{eq:total-f-height}
 \sum_{v\in M_k}n_v|\log|f(t)|_v|
 =2[k:\Q]h(f(t))
 =2d[k:\Q]h(t)+O_f(1).
\end{equation}
At a place \(v\in T_0\), if \(|f(t)|_v<1\), then
\(-\log|f(t)|_v\), up to a bounded local term, is the proximity to a zero
of \(f\); if \(|f(t)|_v>1\), then \(\log|f(t)|_v\) is similarly the
proximity to a pole.  Simplicity ensures that each target appears with
multiplicity one.  Pass to one fixed finite extension containing all zeros
and poles.  At each place over \(T_0\), select a nearest zero or pole for
the point under consideration.  There are only finitely many possible
selections, so Theorem~\ref{input:roth} may be applied separately to each
selection and the resulting inequality divided by the extension degree.
Outside a finite exceptional set,
\begin{equation}\label{eq:roth-near-S}
 \sum_{v\in T_0}n_v|\log|f(t)|_v|
 \le(2+\epsilon)[k:\Q]h(t)+O_{f,S,\epsilon}(1).
\end{equation}
The sum over the finite places \(v\notin S\) in
\eqref{eq:total-f-height} is exactly
\(\log\Norm\mathfrak m_{f,S}(t)\).  Subtract
\eqref{eq:roth-near-S} from \eqref{eq:total-f-height} and, if necessary,
replace \(\epsilon\) by a smaller positive number.
\end{proof}

Fix \(N\ge3\), and put
\begin{equation}\label{eq:PN-QN-FN}
 P_N(T)=q_{N-1}(T),\qquad Q_N(T)=q_N(T),\qquad
 F_N(T)=P_N(T)Q_N(T).
\end{equation}
The rational function \(P_N/Q_N\) has degree \(N-1\).  Its finite zeros
and poles are simple and disjoint, and infinity is one additional simple
zero.  Indeed, \(q_j(T)=U_{j-1}(T/2)\) has distinct roots in characteristic
zero, and \eqref{eq:cassini} makes consecutive terms coprime.

\begin{corollary}[two consecutive recurrence terms]
\label{cor:roth-recurrence-pair}
Let \(k,S,N\) be fixed, with \(N\ge3\), and suppose that \(S\) contains the
bad-reduction places of \(P_N/Q_N\).  For every \(\epsilon>0\), all but
finitely many \(t\in\cO_{k,S}\) satisfy
\begin{equation}\label{eq:roth-recurrence-pair}
 \log\Norm\bigl(F_N(t)\cO_{k,S}\bigr)
 \ge(2N-4-\epsilon)\cJ_k(t)-O_{N,S,\epsilon}(1).
\end{equation}
\end{corollary}

\begin{proof}
Cassini's identity shows that \(P_N(t)\) and \(Q_N(t)\) generate the unit
ideal of \(\cO_{k,S}\).  Therefore the reduced numerator--denominator ideal of
\(P_N(t)/Q_N(t)\) is precisely
\(F_N(t)\cO_{k,S}\).  Lemma~\ref{lem:roth-reduced-divisor}, with
\(d=N-1\), gives the coefficient \(2N-4-\epsilon\).  Finally use
Lemma~\ref{lem:J-height-comparison}; its bounded error is absorbed into the
constant because \(k,N,S,\epsilon\) are fixed.
\end{proof}

\subsection{Local density of square divisibility}

The selection argument in Lemma~\ref{lem:moving-bounded-exponents} requires
the following codimension-two estimate modulo \(\mathfrak p^2\).

\begin{lemma}[square-divisor density]
\label{lem:square-divisor-density}
Fix \(N\ge3\).  Outside a finite set of rational primes depending on
\(N\), \(k_y\), and the integral model, the following holds.  The excluded
set contains the primes ramified in \(k_y\), the bad and quaternion-ramified
primes of the group model, the primes below \(2N(N-1)\), the primes dividing
the discriminant of \(F_N\), and the exceptional primes in
Theorem~\ref{input:strong-approximation}.  If
\(p\) is a remaining prime, then in the reduction image modulo \(p^2\),
\begin{equation}\label{eq:square-divisor-density}
 \frac{\#\{g\in Q_{p^2}:F_N(\Tr g)\equiv0\pmod{\mathfrak p^2}
                 \text{ for some }\mathfrak p\mid p\}}
      {\#Q_{p^2}}
 \ll_N p^{-2}.
\end{equation}
\end{lemma}

\begin{proof}
Put \(A=\cO_{k_y}/\mathfrak p^2\), where
\(q=\Norm\mathfrak p\).  The stated exclusions make the reduction of
\(F_N\) square-free.  Hensel's lemma therefore gives
\[
 \#\{t\in A:F_N(t)=0\}=O_N(1).
\]
For one fixed \(t\in A\), write a matrix of trace \(t\) as
\(\left(\begin{smallmatrix}u&b\\c&t-u\end{smallmatrix}\right)\).  Its
determinant is one exactly when
\[
 bc=u(t-u)-1.
\]
There are \(q^2\) choices for \(u\).  Splitting \(b\) according to its
valuation \(0\), \(1\), or \(2\) gives at most \(3q^2\) choices for \((b,c)\)
for each \(u\).  Thus a trace fiber contains at most \(3q^4\) elements,
whereas
\[
 \#\SL_2(A)=q^6(1-q^{-2}).
\]
It follows that the locus \(F_N(\Tr g)=0\pmod{\mathfrak p^2}\) has relative
density \(O_N(q^{-2})\).  Cassini's identity ensures that the two factors
of \(F_N\) have no common root.  There are at most \([k_y:\Q]\) primes
above \(p\), and each has norm at least \(p\).  The union bound gives
\(O_{N,k_y}(p^{-2})\) in the full product of the local groups.
Theorem~\ref{input:strong-approximation} changes this by at most its fixed
index, proving
\eqref{eq:square-divisor-density}.
\end{proof}

\begin{lemma}[height growth in matrix words]
\label{lem:matrix-word-height}
Let \(L\) be a number field, let \(S\) contain the archimedean places, and
let \(\Omega\subset\GL_d(\cO_{L,S})\) be a fixed finite symmetric set.  There
is a constant \(C_\Omega\) such that every entry of a word \(g\) of length
at most \(n\) in \(\Omega\) has height at most
\(C_\Omega n+C_\Omega\).  If \(P\in L[T]\) is fixed and \(S\) has been
enlarged to contain the denominators of its coefficients, then, whenever
\(P(\Tr g)\ne0\),
\begin{equation}\label{eq:matrix-word-principal-ideal}
 \log\Norm\bigl(P(\Tr g)\cO_{L,S}\bigr)
 \le C_{\Omega,P}n+C_{\Omega,P}.
\end{equation}
\end{lemma}

\begin{proof}
For a place \(v\), put
\[
 \|M\|_v^+=\max\{1,|M_{ij}|_v:1\le i,j\le d\}.
\]
Matrix multiplication gives
\[
 \|AB\|_v^+\le c_v\|A\|_v^+\|B\|_v^+,
\]
where \(c_v=1\) at a nonarchimedean place and \(c_v=d\) at an
archimedean place.  Iterating this inequality for a word of length \(n\),
taking logarithms, multiplying by local degrees, and summing over all
places shows that the projective height of its matrix entries is \(O(n)\).
Only finitely many places contribute because the generators and their
inverses are fixed and \(S\)-integral.  The height of each entry, and then
of their sum \(\Tr g\), is therefore \(O(n)\).

Applying the same local maximum estimate to the finitely many monomials of
\(P\) gives the polynomial height inequality
\[
 h(P(\Tr g))\le(\deg P)h(\Tr g)+O_P(1)=O_{\Omega,P}(n).
\]
The value is \(S\)-integral.  By the product formula, the logarithmic norm
of its prime-to-\(S\) principal ideal is no larger than the sum of all
negative local logarithms, which equals the sum of all positive local
logarithms and is \([L:\Q]h(P(\Tr g))\).  This proves
\eqref{eq:matrix-word-principal-ideal}.
\end{proof}

\begin{lemma}[selection of short perturbations with bounded valuations]
\label{lem:moving-bounded-exponents}
Retain the hypotheses of Theorem~\ref{thm:local-expansion}.  Put
\(\Delta=\rho_y(\Pi^{(2)})\), and fix a nonidentity \(b\in\Pi^{(2)}\), an
integer \(N\ge3\), and a sufficiently small \(\theta>0\).  There are a
rationally saturated finite set \(S\) of places of \(k_y\), an integer
\(D\ge1\), and, for all sufficiently large \(m\), elements
\begin{equation}\label{eq:moving-am}
 a_m=b^m h_m,\qquad |h_m|\le\theta m,
\end{equation}
such that, with \(s_m=\Tr_y(a_m)\),
\begin{align}
 \cJ_{k_y}(s_m)
 &\ge m\cJ_{k_y}(\Tr_y b)-O_{b,y}(\theta m)-O_{b,y}(1),
 \label{eq:moving-J}\\
 d_x(o,a_mo)
 &\le m\ell_x(b)+O_{b,x}(\theta m)+O_{b,x}(1),
 \label{eq:moving-displacement}
\end{align}
and the nonzero ideal
\begin{equation}\label{eq:moving-ideal}
 \mathfrak c_m=F_N(s_m)\cO_{k_y,S}
\end{equation}
has every prime exponent at most \(D\) and satisfies
\begin{equation}\label{eq:moving-ideal-height}
 \log\Norm\mathfrak c_m=O_{N,b,y}(m).
\end{equation}
There is a number \(\theta_0>0\), depending on \(b,y\), such that the
constants implicit in \eqref{eq:moving-J},
\eqref{eq:moving-displacement}, and
\eqref{eq:moving-ideal-height} may be chosen uniformly for
\(0<\theta\le\theta_0\).  The set \(S\) and the exponent \(D\) are allowed
to depend on \(\theta\).
\end{lemma}

\begin{proof}
We specify the constants in the order needed for the final limiting
argument.  Enlarge a rationally saturated set \(S_0\) so that
the representation is integral, the restriction-of-scalars model has good
reduction outside \(S_0\), every rational prime ramified in \(k_y\) is
removed, all primes below \(2N(N-1)\) are removed, and the bad-reduction
places of \(P_N/Q_N\) are included.
Let \(h\) be the position after \(L=\lfloor\theta m\rfloor\) steps of the
lazy random walk defined before
Theorem~\ref{input:prime-square-mixing}, and consider \(a=b^mh\).

Choose one fixed finite Galois extension \(E/\Q\), containing \(k_y\),
which splits the quaternion algebra and the characteristic polynomial of
\(b\).  Only
finitely many places \(w\) of \(E\) satisfy
\(\chi_w(\Tr_y b)>0\).  At every such place choose one of the two global
eigenvalues, denoted \(\lambda_w\), so that \(|\lambda_w|_w>1\).  The
two eigenprojectors are
\[
 E_{w,+}=\frac{b-\lambda_w^{-1}I}{\lambda_w-\lambda_w^{-1}},
 \qquad
 E_{w,-}=\frac{\lambda_w I-b}{\lambda_w-\lambda_w^{-1}}.
\]
They satisfy
\[
 E_{w,+}+E_{w,-}=I,
 \qquad
 b^m=\lambda_w^mE_{w,+}+\lambda_w^{-m}E_{w,-}.
\]
Define the fixed regular linear functions
\[
 L_{w,+}(h)=\Tr(E_{w,+}h),\qquad
 L_{w,-}(h)=\Tr(E_{w,-}h).
\]
Since each projector has rank one, \(L_{w,+}(I)=L_{w,-}(I)=1\).  Taking the
trace after multiplying the projector decomposition by \(h\) gives
\begin{equation}\label{eq:eigen-coefficients}
 \Tr_y(b^mh)=\lambda_w^mL_{w,+}(h)
              +\lambda_w^{-m}L_{w,-}(h).
\end{equation}
For each \(w\), take the product of all Galois conjugates over \(\Q\) of
the regular function \(L_{w,+}\).  Let \(\Phi\) be the product of these
finitely many functions, multiplied by one common nonzero denominator.
Each factor is nonzero because it takes the value one at the identity.
Consequently \(\Phi\in\Q[\mathbf G]\setminus\{0\}\), and
\(\Phi(h)\ne0\) implies \(L_{w,+}(h)\ne0\) for every expanding place
\(w\).

By Theorem~\ref{input:lang-weil}, we may choose a sufficiently large prime
\(r\) of good reduction for which the zero set of \(\Phi\) occupies less
than one eighth of the congruence image.  Use the reduction modulo \(r\)
associated with
the model over the current ring \(\cO_{k_y,S_0}\).  Mixing on this one fixed
quotient shows, for all sufficiently large \(m\), that
\begin{equation}\label{eq:phi-probability}
 \Pr(\Phi(h)=0)<\frac14.
\end{equation}
After obtaining this estimate, enlarge \(S_0\) by all places above \(r\).
The enlargement removes the valuations above \(r\) from the divisor ideals
but does not alter the fixed finite quotient in which the walk has already
been evaluated.
Suppose \(\Phi(h)\ne0\).  Then every \(L_{w,+}(h)\) is nonzero.
Lemma~\ref{lem:matrix-word-height} gives height \(O(L)\) for the matrix
entries of \(h\).  Since the coefficients of the projectors are fixed after
\(b\) and the splitting field have been fixed, the height inequalities
\(h(\alpha\beta)\le h(\alpha)+h(\beta)\) and
\(h(\alpha+\beta)\le h(\alpha)+h(\beta)+O(1)\) give
\begin{equation}\label{eq:eigen-coefficient-height}
 h(L_{w,+}(h)),\ h(L_{w,-}(h))=O_{b,y,E}(L).
\end{equation}
For a nonzero algebraic number \(\alpha\), the product formula and
\(h(\alpha^{-1})=h(\alpha)\) imply, at each fixed place \(w\),
\[
 -\log|\alpha|_w\le C_w h(\alpha),
 \qquad \log^+|\alpha|_w\le C_w h(\alpha),
\]
where the constant only reflects the fixed local-degree normalization.
Applying this to \eqref{eq:eigen-coefficient-height} gives
\begin{equation}\label{eq:eigen-coefficient-local}
 -\log|L_{w,+}(h)|_w=O_{b,y,E}(L),
 \qquad
 \log^+|L_{w,-}(h)|_w=O_{b,y,E}(L).
\end{equation}

Now \(\log|\lambda_w|_w=\chi_w(\Tr_y b)>0\).  Let \(c_0\) be the minimum
of these finitely many positive numbers, and let \(C_0\) dominate the
constants in \eqref{eq:eigen-coefficient-local}.  Choose \(\theta>0\) so
small that \(2C_0\theta<c_0\).  This inequality defines an upper bound
\(\theta_0\) for \(\theta\), and the constants remain uniform for
\(0<\theta\le\theta_0\).  Since \(L\le\theta m\), the logarithm of
the absolute value of the first term in
\eqref{eq:eigen-coefficients} is at least
\[
 m\chi_w(\Tr_y b)-C_0L,
\]
whereas that of the second is at most
\[
 -m\chi_w(\Tr_y b)+C_0L.
\]
The first term therefore dominates exponentially.  At a nonarchimedean
place the ultrametric inequality makes the norm of the sum equal to the norm
of the first term.  At an archimedean place the ordinary triangle inequality
loses only \(O(1)\) in the logarithm.  Thus
\[
 \log|\Tr_y(b^mh)|_w
 \ge m\chi_w(\Tr_y b)-O_{b,y,E}(L)-O_{b,y,E}(1).
\]
At a nonarchimedean place with the displayed trace norm greater than one,
\(\chi_v(t)=\log|t|_v\); at an archimedean place the two quantities differ
by at most \(\log2\), as in the proof of
Lemma~\ref{lem:J-height-comparison}.  Hence
\[
 \chi_w(\Tr_y(b^mh))
 \ge m\chi_w(\Tr_y b)-O_{b,y,E}(L)-O_{b,y,E}(1).
\]
Summing with the local-degree weights over the finitely many expanding
places gives
\[
 \cJ_E(\Tr_y(b^mh))
 \ge m\cJ_E(\Tr_yb)-O_{b,y,E}(\theta m)-O_{b,y,E}(1);
\]
the remaining local summands on the left are nonnegative.  Both traces lie
in \(k_y\), so the base-change formula
\eqref{eq:J-base-change} says
\[
 \cJ_E(\Tr_y(b^mh))=[E:k_y]\cJ_{k_y}(\Tr_y(b^mh)),
 \quad
 \cJ_E(\Tr_yb)=[E:k_y]\cJ_{k_y}(\Tr_yb).
\]
Division by \([E:k_y]\) proves \eqref{eq:moving-J} with the field
normalization stated in the lemma.

For the displacement estimate, retain the basepoint \(o\) fixed earlier and
choose a point \(z\) on the \(x\)-axis of \(b\).  The distance
\(d_x(o,z)\) is a fixed constant depending only on \(b,x,o\).  Then
\[
 d_x(o,b^m h o)
 \le d_x(o,z)+d_x(z,b^mz)+d_x(b^mz,b^mho)
 \le m\ell_x(b)+2d_x(o,z)+d_x(o,ho).
\]
The displacement of a word is at most the sum of the displacements of its
letters.  Since the generating set is fixed and \(|h|\le\theta m\), this
proves \eqref{eq:moving-displacement}.

The word length of \(b^mh\) is \(O(m)\).  Apply
Lemma~\ref{lem:matrix-word-height} with \(P=F_N\).  After enlarging
\(S_0\) by the fixed coefficient denominators, it gives a constant \(A_N\),
uniform for \(0<\theta\le\theta_0\), such that, whenever the value is
nonzero,
\begin{equation}\label{eq:FN-height-bound}
 \log\Norm\bigl(F_N(\Tr_y(b^mh))\cO_{k_y,S_0}\bigr)
 \le A_Nm+O_N(1).
\end{equation}
The distinguished real place makes \(\cJ_{k_y}(\Tr_y b)>0\), so
\eqref{eq:moving-J} implies that the traces under consideration eventually
leave every fixed finite set.  In particular, \(F_N\) does not vanish for
all sufficiently large \(m\) on the event \(\Phi(h)\ne0\).

Choose a large constant \(Z\), larger than every exceptional rational prime
in Lemma~\ref{lem:square-divisor-density} and
Theorem~\ref{input:prime-square-mixing}, and enlarge \(S_0\) by every place
above a rational prime at most \(Z\); call the result \(S\).  By
Lemma~\ref{lem:square-divisor-density} and
Theorem~\ref{input:prime-square-mixing}, for a good rational prime \(p>Z\),
\begin{equation}\label{eq:square-event}
 \Pr\bigl(\mathfrak p^2\mid F_N(\Tr_y(b^mh))
       \text{ for some }\mathfrak p\mid p\bigr)
 \le C_Np^{-2}+p^B e^{-\delta L}.
\end{equation}
Theorem~\ref{input:prime-square-mixing} is uniform in the starting point,
so it applies with starting point \(b^m\).

If some exponent in \(F_N(\Tr_y(b^mh))\cO_{k_y,S}\) is greater than an
integer \(D\), then \eqref{eq:FN-height-bound} shows that the rational prime
below it is at most
\[
 Y_m=\exp\left(\frac{A_Nm+O_N(1)}{D+1}\right).
\]
Indeed, \((D+1)\log\Norm\mathfrak p\le A_Nm+O_N(1)\), and the rational
prime below \(\mathfrak p\) is no larger than \(\Norm\mathfrak p\).  The
union bound in \eqref{eq:square-event} gives
\begin{align}
 \Pr(&\text{some prime-to-}S\text{ exponent exceeds }D)\notag\\
 &\le C_N\sum_{p>Z}p^{-2}
 +\exp\left(
   \left(\frac{(B+1)A_N}{D+1}-\delta\theta\right)m+o(m)
          \right).
 \label{eq:bounded-exponent-union}
\end{align}
Choose \(Z\) first so that the first term is less than \(1/8\).  Then choose
the fixed integer \(D\) so that
\begin{equation}\label{eq:explicit-sieve-exponent-choice}
 D+1>\frac{(B+1)A_N}{\delta\theta}.
\end{equation}
This makes the exponent in the second term negative.  Here \(B,\delta\)
come from the fixed prime-square mixing estimate, \(A_N\) was fixed by the
height bound after \(N,b,y\) were fixed, and \(\theta>0\) was fixed before
\(D\).  Thus \(D\) is independent of \(m\), although it may depend on
\(N,b,y,\theta\), as allowed in the statement of the lemma.
For large \(m\), \eqref{eq:bounded-exponent-union} is less than \(1/4\).
Together with \eqref{eq:phi-probability}, the total probability of the
excluded events is less than one.  Hence some walk position \(h_m\)
satisfies all the stated properties.  Equation
\eqref{eq:moving-ideal-height} follows from
\eqref{eq:FN-height-bound}.
\end{proof}

\section{Controlling fibers of two trace functions}
\label{sec:two-traces}

An edge in the trace graph is the pair
\((\Tr_y(AH),\Tr_y H)\).  To recover enough distinct edges from a large
set of group elements, we must show that no such pair is assumed
exponentially often.  The required bound is ultimately a divisor estimate
over a number field.

\begin{lemma}[quadratic norm divisor estimate]
\label{lem:quadratic-norm-divisor}
Let \(L\) be a number field and \(S\) a finite set of places containing the
archimedean places.  Fix \(A>0\).
Uniformly in \(d,N\in\cO_{L,S}\) with \(dN\ne0\) and
\(h(d),h(N)\le Am\), the number of pairs
\(u,v\in\cO_{L,S}\) satisfying
\[
 h(u),h(v)\le Am,
 \qquad u^2-dv^2=N
\]
is at most \(\exp(o_{A,L,S}(m))\).
\end{lemma}

\begin{proof}
Suppose first that \(d\) is not a square in \(L\), and put
\[
 E=L(T),\qquad T^2=d.
\]
This is a quadratic field extension of \(L\).  Let \(S_E\) be the places
of \(E\) above \(S\), and let \(\cO_{E,S_E}\) be the integral closure of
\(\cO_{L,S}\).  Since \(T\) is integral over \(\cO_{L,S}\), every solution
gives
\[
 \alpha=u+vT\in\cO_{E,S_E},\qquad
 \alpha\bar\alpha=N.
\]
The map \((u,v)\mapsto\alpha\) is injective because \(1,T\) is an
\(L\)-basis of \(E\).

We first count the possible prime-to-\(S_E\) principal ideals
\((\alpha)\).  For a prime \(\mathfrak P\notin S_E\), integrality of both
\(\alpha\) and \(\bar\alpha\) gives
\[
 0\le v_{\mathfrak P}(\alpha)
 \le v_{\mathfrak P}(N).
\]
If \(\mathfrak p\) is the prime of \(L\) below \(\mathfrak P\), then
\(v_{\mathfrak P}(N)\le2v_{\mathfrak p}(N)\), and at most two primes of
\(E\) lie above \(\mathfrak p\).  The number of possible ideal patterns is
therefore at most
\begin{equation}\label{eq:quadratic-ideal-patterns}
 \prod_{\substack{\mathfrak p\notin S\\
                   v_{\mathfrak p}(N)>0}}
       \bigl(2v_{\mathfrak p}(N)+1\bigr)^2.
\end{equation}
The logarithmic norm of the prime-to-\(S\) numerator ideal of \(N\) is
\(O_{L,S,A}(m)\).  Indeed, its valuation sum is the finite-place
contribution to the height of \(N^{-1}\), and
\(h(N^{-1})=h(N)\le Am\).

We claim that \eqref{eq:quadratic-ideal-patterns} is
\(\exp(o(m))\) uniformly in \(d,N\).  To see this, fix \(T_0>1\).  Only
finitely many prime ideals of \(L\) have norm at most \(T_0\), and their
contribution to
the logarithm of \eqref{eq:quadratic-ideal-patterns} is
\(O_{L,T_0}(\log m)\).  For the remaining primes use
\(\log(2e+1)\le e\log3\) and
\[
 \sum_{\Norm\mathfrak p>T_0}v_{\mathfrak p}(N)
 \le\frac1{\log T_0}
     \sum_{\mathfrak p}v_{\mathfrak p}(N)
                           \log\Norm\mathfrak p
 =O(m/\log T_0).
\]
First let \(m\to\infty\), and then \(T_0\to\infty\).  This proves the
required uniform subexponential estimate for the possible principal ideals.

It remains to count generators of each possible principal ideal.  If
\(\alpha\) and \(\alpha'\) arise from two admissible solutions with the same prime-to-
\(S_E\) principal ideal, then
\[
 \varepsilon=\alpha/\alpha'\in\cO_{E,S_E}^{\times}.
\]
The height inequalities for sums, products, and quotients, together with
\(h(T)=\tfrac12h(d)\), give
\[
 h(\alpha),h(\alpha')=O_{L,S,A}(m),\qquad
 h(\varepsilon)=O_{L,S,A}(m).
\]
For an \(S_E\)-unit define its weighted logarithmic vector by
\[
 \ell_E(\varepsilon)
   =\bigl(n_w\log|\varepsilon|_w\bigr)_{w\in S_E},
\]
where \(n_w\) is the usual local-degree weight.  The product formula and
the fact that \(|\varepsilon|_w=1\) outside \(S_E\) give the exact identity
\begin{equation}\label{eq:unit-log-height-identity}
 \|\ell_E(\varepsilon)\|_1
 =2[E:\Q]h(\varepsilon).
\end{equation}
Every such field \(E\) has degree at most \(2[L:\Q]\), and \(S_E\) has
cardinality bounded in terms of \(L,S\).  Kronecker's theorem says that
height zero is equivalent to being a root of unity.  Together with
Northcott finiteness, it follows that there is a positive lower bound,
depending only on \(2[L:\Q]\), for the height of a non-torsion algebraic
number of these degrees; the number of roots of unity is uniformly bounded
as well \cite[Sections~1.5--1.6]{BombieriGubler}.  By
\eqref{eq:unit-log-height-identity}, distinct logarithmic vectors modulo
torsion are therefore separated by a fixed positive distance.  They lie in
a ball of radius \(O(m)\) in a space of uniformly bounded dimension.
Packing gives
\[
 \#\{\varepsilon\in\cO_{E,S_E}^{\times}:h(\varepsilon)=O(m)\}
 =m^{O_{L,S,A}(1)}.
\]
Thus each possible principal ideal contributes only polynomially many
solutions, and
the nonsplit case follows.

It remains to treat the split case explicitly.  Write \(d=a^2\) with
\(a\in L^\times\).  Since \(d\in\cO_{L,S}\), valuations outside \(S\)
show that \(a\in\cO_{L,S}\), and \(h(a)=\tfrac12h(d)\).  Every solution
gives two nonzero \(S\)-integers
\[
 \alpha_+=u+av,\qquad \alpha_-=u-av,\qquad
 \alpha_+\alpha_-=N.
\]
The height bounds on \(u,v,a\) give
\(h(\alpha_+),h(\alpha_-)=O_{L,S,A}(m)\).  The prime-to-\(S\) ideal of
\(\alpha_+\) divides that of \(N\), so the
same small-prime/large-prime argument gives \(\exp(o(m))\) possible ideals.
For each fixed ideal, two possible generators differ by an \(S\)-unit of
height \(O(m)\).  The logarithmic-vector argument above, now in the fixed
field \(L\), gives only \(m^{O_{L,S,A}(1)}\) generators.  Finally,
\(\alpha_+\) determines \(\alpha_-=N/\alpha_+\), and then determines
\(u=(\alpha_++\alpha_-)/2\) and
\(v=(\alpha_+-\alpha_-)/(2a)\).  Hence the map from solutions to
\(\alpha_+\) is injective.  This proves the same subexponential bound in
the split case and completes the proof.
\end{proof}

\begin{lemma}[uniform fibers of two trace functions]
\label{lem:moving-two-trace-fibers}
Retain \(x,y,\Delta\) from Theorem~\ref{thm:local-expansion}.  Suppose that
\(A_m\in\Delta\) is noncentral and has word length at most \(Wm\), where
\(W\) is fixed.  For every fixed \(C>0\), uniformly in
\(\xi,\zeta\in k_y\),
\begin{equation}\label{eq:moving-two-trace-fiber}
 \#\{H\in\Delta:d_x(o,Ho)\le Cm,
       \ \Tr_yH=\xi,\ \Tr_y(A_mH)=\zeta\}
 \le\exp(o_{C,W,x,y}(m)).
\end{equation}
The same conclusion holds for squared traces, with at most four choices of
sign.
\end{lemma}

\begin{proof}
Fix a splitting field of the quaternion algebra.  Word length
and \(x\)-displacement are comparable, so all matrix entries of \(A_m\) and
of every \(H\) counted in \eqref{eq:moving-two-trace-fiber} have logarithmic
height \(O_{C,W,x,y}(m)\) by Lemma~\ref{lem:matrix-word-height}, and they
lie in one fixed ring of \(S\)-integers.
If the fiber is empty, there is nothing to prove.  If it is nonempty, choose
one matrix \(H_0\) in it.  Then
\[
 \xi=\Tr_yH_0,
 \qquad \zeta=\Tr_y(A_mH_0)
\]
also have height \(O_{C,W,x,y}(m)\).  Thus, whenever the fiber is nonempty,
its prescribed trace values \(\xi,\zeta\) satisfy the same height bound as
its matrices.  All constants below can therefore be chosen independently
of \(\xi,\zeta\).

Write
\[
 A=\begin{pmatrix}\alpha&\beta\\ \gamma&\delta\end{pmatrix},
 \qquad
 H=\begin{pmatrix}r&s\\ u&\xi-r\end{pmatrix},
 \qquad w=\zeta-\delta\xi,
\]
and put
\[
 \Delta_A=(\alpha+\delta)^2-4,\qquad
 y_0=2\beta r-(\alpha-\delta)s,
 \qquad \kappa=\beta\{\xi(\alpha-\delta)-2w\}.
\]
Since \(A\) is noncentral in a faithful cocompact Fuchsian group,
\(\Tr(A)^2\ne4\): an element with trace squared equal to four is either
central or projects to a parabolic element, and the latter cannot occur in
a cocompact lattice.  Hence \(\Delta_A\ne0\).  Suppose first that
\(\beta\ne0\).  The determinant and
second trace equations are, respectively,
\begin{align}
 su&=r(\xi-r)-1,
 \label{eq:fiber-det-scalar}\\
 \beta u&=w-(\alpha-\delta)r-\gamma s.
 \label{eq:fiber-trace-scalar}
\end{align}
Multiply \eqref{eq:fiber-det-scalar} by \(\beta\), substitute
\eqref{eq:fiber-trace-scalar}, and collect the terms in \(r\).  This gives
\[
 \beta r^2-\{\beta\xi+(\alpha-\delta)s\}r
      +\beta+ws-\gamma s^2=0.
\]
Multiplication by \(4\beta\) and completion of the square gives
\[
 (y_0-\beta\xi)^2
 =\Delta_As^2+2\kappa s+\beta^2(\xi^2-4),
\]
because
\(\Delta_A=(\alpha-\delta)^2+4\beta\gamma\).  Multiplying this identity by
\(\Delta_A\) and completing the square once more gives
\begin{equation}\label{eq:quadratic-elimination}
 (\Delta_As+\kappa)^2-\Delta_A(y_0-\beta\xi)^2
 =\kappa^2-\Delta_A\beta^2(\xi^2-4).
\end{equation}
The map
\[
 H\longmapsto
 (Y,V)=(y_0-\beta\xi,\Delta_As+\kappa)
\]
is injective: \(V\) recovers \(s\), then \(Y\) recovers \(r\), and the
second trace equation recovers \(u\).  All coefficients and solutions in
\eqref{eq:quadratic-elimination} have height \(O_{C,W,x,y}(m)\).

The right side of \eqref{eq:quadratic-elimination} equals
\begin{equation}\label{eq:fricke-degeneracy}
 4\beta^2\bigl(\Tr[A,H]-2\bigr),
\end{equation}
where \([A,H]=AHA^{-1}H^{-1}\).  To verify
\eqref{eq:fricke-degeneracy}, use
\begin{equation}\label{eq:fricke-commutator}
 \Tr[A,H]
 =\Tr(A)^2+\Tr(H)^2+\Tr(AH)^2
  -\Tr(A)\Tr(H)\Tr(AH)-2.
\end{equation}
Indeed, substitution of \(\Tr(A)=\alpha+\delta\),
\(\Tr(H)=\xi\), \(\Tr(AH)=\zeta\), and
\(w=\zeta-\delta\xi\) gives
\[
 \kappa^2-\Delta_A\beta^2(\xi^2-4)
 =4\beta^2\{\Tr(A)^2+\xi^2+\zeta^2
       -\Tr(A)\xi\zeta-4\},
\]
which is \eqref{eq:fricke-degeneracy} by
\eqref{eq:fricke-commutator}.

If the right side is nonzero,
put
\[
 d=\Delta_A,\qquad
 N_A=\kappa^2-\Delta_A\beta^2(\xi^2-4),\qquad
 Y=y_0-\beta\xi,\qquad V=\Delta_As+\kappa.
\]
Then \eqref{eq:quadratic-elimination} is exactly
\(V^2-dY^2=N_A\).  The entry-height estimates above and the elementary
height inequalities for sums and products give
\[
 h(d),\ h(N_A),\ h(Y),\ h(V)=O_{C,W,x,y}(m).
\]
Thus all hypotheses of Lemma~\ref{lem:quadratic-norm-divisor} hold over the
fixed base splitting field and its fixed ring of \(S\)-integers; the lemma is
uniform in the varying coefficient \(d\) and its associated quadratic
algebra.  It gives
\(\exp(o_{C,W,x,y}(m))\) possible
pairs \((Y,V)\), hence the same number of matrices \(H\).

If \(N_A=0\), equivalently, if the right-hand side of
\eqref{eq:quadratic-elimination} is zero, then \(\Tr[A,H]=2\).  Since
\([A,H]\in\Delta\), its projective image is
either the identity or a parabolic element.  The latter cannot occur in a
cocompact Fuchsian group, so \([A,H]=I\).  The centralizer of a nonidentity surface-group
element is cyclic.  If its primitive hyperbolic generator has length
\(\ell_0>0\), the absolute trace of its \(n\)-th power is
\(2\cosh(|n|\ell_0/2)\).  A prescribed signed trace therefore permits at
most two exponents \(n\), up to the fixed central choice of lift.  Hence the
degenerate fiber has \(O(1)\) elements.

If \(\beta=0\) and \(\gamma\ne0\), replace \(A,H\) by their transposes.
Transposition preserves determinants and each of the three traces
\(\Tr A\), \(\Tr H\), and \(\Tr(AH)\), while the upper-right entry of
\(A^{\mathsf T}\) is \(\gamma\ne0\).  The preceding argument therefore
applies without changing the fiber.  If
\(\beta=\gamma=0\), the two trace equations determine both diagonal entries
of \(H\), while \(\det H=1\) fixes the product \(su\).  A nonzero product
has \(\exp(o(m))\) factorizations.  Indeed, if \(su=c\ne0\), then
\((s)\) divides \((c)\) outside \(S\).  The number of possible divisor
ideals is bounded by
\[
 \prod_{\mathfrak p\notin S}
       \bigl(v_{\mathfrak p}(c)+1\bigr)=\exp(o(m)),
\]
by the small-prime/large-prime argument in
\eqref{eq:quadratic-ideal-patterns}, since \(h(c)=O(m)\).  For a fixed
principal ideal, two possible generators differ by an \(S\)-unit; the
bounded-degree height-separation and packing argument from
Lemma~\ref{lem:quadratic-norm-divisor} gives only \(m^{O(1)}\) such
generators of height \(O(m)\).  Once \(s\) is known, \(u=c/s\) is fixed.
Thus the number of factorizations is \(\exp(o(m))\).  If \(su=0\), the
Fricke identity gives \(\Tr[A,H]=2\), so the commutator is either the
identity or a
nontrivial parabolic element.  The latter cannot occur in a cocompact
Fuchsian group; hence \([A,H]=I\).  Since \(A\) is noncentral diagonal with
distinct eigenvalues, a matrix commuting with \(A\) is diagonal, and
therefore \(s=u=0\).  This leaves only the \(O(1)\) diagonal solutions with
the prescribed trace.  Finally, two squared-trace conditions split into at
most four signed conditions.
\end{proof}

\section{Congruence counting in the trace graph}
\label{sec:congruence-wedge}

The following combinatorial lemma counts ordered length-two paths whose two
edges share their middle vertex.  It does not assume that the bipartite
graph contains a Cartesian product.

\begin{lemma}[congruence count for length-two paths]
\label{lem:congruence-wedge}
Let \(\mathcal R\) be a Dedekind domain, let
\(P,Q\in\mathcal R\setminus\{0\}\) generate
coprime ideals, and let \(\varnothing\ne G\subset U\times V\), where
\(U,V\subset\mathcal R\) are finite.  Put
\[
 X_- =\{Qu-Pv:(u,v)\in G\},\qquad
 X_+ =\{Qv-Pu:(u,v)\in G\},
\]
and
\[
 M_{(PQ)}(V)=
 \max_{c\in\mathcal R/(PQ)}\#\bigl(V\cap(c+(PQ))\bigr).
\]
Here \((PQ)\) is the principal ideal generated by \(PQ\), and
\(V\cap(c+(PQ))\) denotes the elements of \(V\) whose residue class in
\(\mathcal R/(PQ)\) is \(c\).
The nonemptiness of \(G\) implies that \(V\) is nonempty and
\(M_{(PQ)}(V)\ge1\), so the denominators below are nonzero.
Then
\begin{equation}\label{eq:congruence-wedge}
 \#X_-\,\#X_+
 \ge \frac1{M_{(PQ)}(V)}\sum_{v\in V}\deg_G(v)^2
 \ge \frac{(\#G)^2}{\#V\,M_{(PQ)}(V)}.
\end{equation}
\end{lemma}

\begin{proof}
An ordered length-two path is a triple \((u,v,u')\) with
\((u,v),(u',v)\in G\).  Send it to
\[
 (Qu-Pv,\ Qv-Pu')\in X_-\times X_+.
\]
For a fixed image \((x,z)\), reduction of the two equations gives
\[
 v\equiv-P^{-1}x\pmod{(Q)},\qquad
 v\equiv Q^{-1}z\pmod{(P)}.
\]
The inverses exist because \((P)+(Q)=\mathcal R\).  The Chinese remainder theorem
therefore confines \(v\) to one residue class modulo \((PQ)\).  Once
\(v,x,z\) are fixed, the two equations determine \(u,u'\) in the fraction
field.  Hence every fiber of this map has cardinality at most
\(M_{(PQ)}(V)\).  There are \(\sum_v\deg_G(v)^2\) paths.  This proves the
first inequality, and Cauchy--Schwarz gives
\(\sum_v\deg_G(v)^2\ge(\#G)^2/\#V\).
\end{proof}

\begin{theorem}[upper bound for the recurrence divisor]
\label{thm:recurrence-divisor-upper}
Retain \(x,y,\Delta\), assume the signed trace set at \(x\) satisfies the
critical trace-growth bound, and fix \(N\ge3\).  Let
\(A_m\in\Delta\) be noncentral, interpreted in \(\Pi^{(2)}\) through
\(\rho_y^{-1}\), and suppose its word length is at most \(Wm\) for one
fixed \(W\).  Put
\[
 s_m=\Tr_y(A_m),\qquad D_m=d_x(o,A_mo).
\]
Suppose \(D_m=O(m)\).  Let \(S\) and \(D\) be fixed, with \(S\)
rationally saturated, and assume that the nonzero ideals
\[
 \mathfrak c_m=
 \bigl(q_{N-1}(s_m)q_N(s_m)\bigr)\cO_{k_y,S}
\]
have every prime exponent at most \(D\) and satisfy
\begin{equation}\label{eq:upper-ideal-linear-bound}
 \log\Norm\mathfrak c_m\le Am+O(1)
\end{equation}
for one fixed constant \(A\).  Then
\begin{equation}\label{eq:recurrence-divisor-upper}
 \log\Norm\mathfrak c_m
 \le\frac{2N-1}{2}D_m+o(m).
\end{equation}
\end{theorem}

\begin{proof}
Let \(\delta_D\) be the spectral exponent in
Theorem~\ref{input:bounded-exponent-counting}.  Choose
\begin{equation}\label{eq:moving-wedge-radius}
 C>\frac{3A}{\delta_D}
\end{equation}
and put \(R=Cm\).  Equations
\eqref{eq:upper-ideal-linear-bound} and
\eqref{eq:moving-wedge-radius} give, with a linear margin,
\[
 3\log\Norm\mathfrak c_m<\delta_D R-o(R).
\]
Thus the absorption condition
\eqref{eq:counting-absorption-condition} holds.
Theorem~\ref{input:bounded-exponent-counting}, in the form
\eqref{eq:bounded-exponent-flat-form}, gives
\begin{equation}\label{eq:ball-flatness}
 \max_c\#\{g\in\cB_R^x:\Tr_y g\equiv c
                    \pmod{\mathfrak c_m}\}
 \le\frac{e^{R+o(m)}}{\Norm\mathfrak c_m}.
\end{equation}
Here \(D\), \(A\), and \(\delta_D\) are fixed before \(C\) is chosen;
all four constants remain fixed as \(m\to\infty\).

Form a simple bipartite graph whose left and right vertices are the trace
values that occur in the first and second coordinates, respectively, and
whose edge set is
\begin{equation}\label{eq:moving-trace-graph}
 \cG_m=\{(u,v)=(\Tr_y(A_mg),\Tr_y g):g\in\cB_R^x\}.
\end{equation}
Orbit counting gives \(e^{R+O(1)}\) group elements in the ball, and
Lemma~\ref{lem:moving-two-trace-fibers} bounds every fiber of the two trace
coordinates by \(e^{o(m)}\).  Therefore
\begin{equation}\label{eq:graph-edges}
 \#\cG_m\ge e^{R-o(m)}.
\end{equation}
Let \(V_m\) be the second vertex class.  Persistence of squared-trace
collisions from \(x\) to \(y\), together with the critical trace-growth
bound at \(x\), gives
\begin{equation}\label{eq:V-size}
 \#V_m\le e^{R/2+o(m)}.
\end{equation}
Indeed, the ball has at most \(e^{R/2+o(R)}\) squared-trace values at \(x\).
Collision persistence from Lemma~\ref{lem:collision-persistence} allows no
more squared-trace values at \(y\), and each squared trace has at most two
signed square roots.  Here \(R=Cm\), with \(C\) fixed before
\(m\to\infty\).

Put
\[
 V'_m=\left\{v\in V_m:
       \deg_{\cG_m}(v)\ge\frac{\#\cG_m}{4\#V_m}\right\},
 \qquad
 \cG'_m=\{(u,v)\in\cG_m:v\in V'_m\}.
\]
The vertices outside \(V'_m\) are incident to fewer than
\(\#V_m\cdot\#\cG_m/(4\#V_m)=\#\cG_m/4\) edges, so at most one quarter
of the edges is deleted.  For every
\((u,v)\in\cG'_m\), choose one element
\(g_{u,v}\in\cB_R^x\) satisfying
\[
 \Tr_y(A_mg_{u,v})=u,\qquad \Tr_y(g_{u,v})=v.
\]
For each fixed \(v\in V'_m\), the representatives chosen for its incident
edges are distinct: one group element determines only one ordered pair of
traces.  The definition of \(V'_m\), together with
\eqref{eq:graph-edges} and \eqref{eq:V-size}, therefore gives at least
\(e^{R/2-o(m)}\) distinct chosen group elements with second trace \(v\).
Combining this lower bound with \eqref{eq:ball-flatness} gives
\begin{equation}\label{eq:unweighted-flatness}
 M_{\mathfrak c_m}(V'_m)
 \le\frac{e^{R/2+o(m)}}{\Norm\mathfrak c_m}.
\end{equation}
Indeed, if one residue class contains \(r\) values from \(V'_m\), the
chosen representatives of their incident edges give at least
\(r e^{R/2-o(m)}\) elements of \(\cB_R^x\) whose traces lie in that class.
Equation~\eqref{eq:ball-flatness} bounds their number by
\(e^{R+o(m)}/\Norm\mathfrak c_m\), which yields the displayed inequality.
The passage to \(V'_m\) reconciles two different counts: congruence
counting counts group elements, with all trace multiplicities retained,
whereas Lemma~\ref{lem:congruence-wedge} counts distinct trace values.

Apply Lemma~\ref{lem:congruence-wedge} over \(\cO_{k_y,S}\), with
\[
 P=q_{N-1}(s_m),\qquad Q=q_N(s_m).
\]
Cassini's identity makes \((P)\) and \((Q)\) coprime, and
\(\mathfrak c_m=(PQ)\).  Thus
\(M_{\mathfrak c_m}(V'_m)=M_{(PQ)}(V'_m)\) in the notation of
Lemma~\ref{lem:congruence-wedge}.  Equations
\eqref{eq:graph-edges}--\eqref{eq:unweighted-flatness} give
\begin{equation}\label{eq:wedge-lower}
 \#X_-\,\#X_+
 \ge e^{R-o(m)}\Norm\mathfrak c_m.
\end{equation}

The recurrence identities show that every element of \(X_-\) and \(X_+\)
is a trace of a group element:
\[
 X_- =\{\Tr_y(A_m^Ng_{u,v}):(u,v)\in\cG'_m\},\qquad
 X_+ =\{\Tr_y(A_m^{-(N-1)}g_{u,v}):(u,v)\in\cG'_m\}.
\]
For \((u,v)\in\cG'_m\) and \(j\ge0\),
\[
 d_x(o,A_m^jg_{u,v}o)\le R+jD_m,
\]
and the analogous inequality holds for negative powers.  For each of these
two displacement balls, persistence of squared-trace equalities from \(x\)
to \(y\) bounds the number of traces at \(y\) by the number of squared
traces at \(x\), up to a factor two for the choice of sign.  The critical
trace-growth bound therefore implies
\begin{equation}\label{eq:wedge-upper}
 \#X_-\,\#X_+
 \le\exp\left(R+\frac{2N-1}{2}D_m+o(m)\right).
\end{equation}
The error in the estimate for \(X_-\) is \(o(R+ND_m)\), and that for
\(X_+\) is \(o(R+(N-1)D_m)\).  Both are \(o(m)\), because \(C\) and \(N\)
are fixed at this stage and \(D_m=O(m)\).
Comparison of \eqref{eq:wedge-lower} and \eqref{eq:wedge-upper} proves
\eqref{eq:recurrence-divisor-upper}.
\end{proof}

\begin{remark}[sharpness of the trace-growth exponent for this argument]
\label{rem:critical-growth-sharpness}
Suppose more generally that
\(N_\Gamma(X)\le X^{\sigma+o(1)}\), and put \(\alpha=\sigma/2\).  A ball of
radius \(R\) then has at most \(\exp(\alpha R+o(R))\) trace values.  In the
preceding proof this gives
\[
 \#V_m\le e^{\alpha R+o(m)},\qquad
 \deg_{\cG'_m}(v)\ge e^{(1-\alpha)R-o(m)},\qquad
 M_{\mathfrak c_m}(V'_m)
 \le\frac{e^{\alpha R+o(m)}}{\Norm\mathfrak c_m}.
\]
The two shifted trace supports are bounded at radii \(R+ND_m\) and
\(R+(N-1)D_m\).  Substitution in the length-two-path lower bound and the
shifted-trace upper bound consequently gives
\[
 \log\Norm\mathfrak c_m
 \le (2\sigma-2)R
      +\frac{\sigma(2N-1)}2D_m+o(m).
\]
At \(\sigma=1\), the additional term involving \(R=Cm\) cancels and the
coefficient used in Theorem~\ref{thm:local-expansion} is recovered.  A fixed
power loss \(\sigma>1\) leaves a positive multiple of the auxiliary radius
and does not prove the local theorem.  The hypothesis
\(N_\Gamma(X)\le X^{1+o(1)}\) is the critical case treated by the theorem.
\end{remark}

\section{Proof of the local expansion theorem}
\label{sec:proof-local}

\begin{proof}[Proof of Theorem~\ref{thm:local-expansion}]
For the identity, \(\Tr_y b=2\), every local expansion exponent vanishes,
and \(\ell_x(b)=0\); hence the assertion holds.  Let
\(b\in\Pi^{(2)}\) be
nonidentity.  Its image under every faithful cocompact Fuchsian
representation is hyperbolic in \(\PSL_2\), and hence every
\(\SL_2\)-lift is noncentral.

Fix, in this order, an integer \(N\ge3\), a number
\(0<\epsilon<2N-4\), and a sufficiently small positive number \(\theta\);
then fix the set
\(S\), exponent \(D\), and sequence \(a_m\) given by
Lemma~\ref{lem:moving-bounded-exponents}.  Finally choose the ball constant
\(C\) required in the proof of
Theorem~\ref{thm:recurrence-divisor-upper}.  These data remain fixed as
\(m\to\infty\).

Equation \eqref{eq:moving-J} shows that \(s_m=\Tr_y(a_m)\) leaves every
finite set.  Corollary~\ref{cor:roth-recurrence-pair} therefore applies for
all sufficiently large \(m\).  Combining it with
Theorem~\ref{thm:recurrence-divisor-upper} gives
\begin{align*}
 &(2N-4-\epsilon)
 \bigl(m\cJ_{k_y}(\Tr_y b)-O(\theta m)-O(1)\bigr)\\
 &\hspace{22mm}\le
 \frac{2N-1}{2}
 \bigl(m\ell_x(b)+O(\theta m)+O(1)\bigr)+o(m).
\end{align*}
Divide first by \(m\) and let \(m\to\infty\), with
\(N,\epsilon,\theta,S,D,C\) fixed.  To take the next limit, repeat the
entire construction for a family of positive values \(\theta\to0\).  Each
new value of \(\theta\) is allowed to produce a new set \(S_\theta\), a new
exponent bound \(D_\theta\), a new spectral exponent, and a new ball
constant \(C_\theta\).  For each value of \(\theta\), the limit in \(m\)
is taken before \(\theta\) is changed.  Consequently these auxiliary data
need not be uniform as \(\theta\to0\).  The error coefficients in
\eqref{eq:moving-J} and \eqref{eq:moving-displacement} are uniform for
\(0<\theta\le\theta_0\), so this family of inequalities permits
\(\theta\to0\).  We then let \(\epsilon\to0\), and finally
\(N\to\infty\).  Since
\[
 \lim_{N\to\infty}\frac{(2N-1)/2}{2N-4}=\frac12,
\]
we obtain
\[
 \cJ_{k_y}(\Tr_y b)\le\frac{\ell_x(b)}2.
\]
This is \eqref{eq:main-local}.
\end{proof}

\section{Arithmeticity and exclusion of transcendental characters}
\label{sec:arithmeticity}

\begin{proof}[Proof of Theorem~\ref{thm:main-arithmeticity}]
Let \(x\) be the marked Fuchsian character of \(\Gamma\).

Suppose first that \(x\) is algebraic.  Apply
Theorem~\ref{thm:local-expansion} with \(y=x\), and put
\(\Lambda=\Gamma^{(2)}\).  Since \(\Pi\) is finitely generated and
\(\Pi/\Pi^{(2)}\) is an elementary abelian two-group, \(\Lambda\) has
finite index in \(\Gamma\).  In particular, \(\Lambda\) is itself a
torsion-free cocompact lattice.  Choose an \(\SL_2\)-lift \(\rho_x\) of the
marked representation and set
\(\Delta_x=\rho_x(\Pi^{(2)})\); this is a linear lift of \(\Lambda\).
Its full inverse image is
\(\widetilde\Lambda=\{\pm h:h\in\Delta_x\}\), so the full inverse image and
\(\Delta_x\) generate the same ordinary trace field.
The invariant-trace-field identity of
Lemma~\ref{lem:finite-trace-coordinates} gives
\[
 \Q(\Tr\Delta_x)
 =\Q\bigl(\Tr(\widetilde\gamma)^2:
             \widetilde\gamma\in\widetilde\Gamma\bigr)
 =k_\Gamma=k_x.
\]
Thus the ordinary trace field of \(\Lambda\), computed from its full
inverse image, is the invariant trace field of \(\Gamma\).
For every nonidentity
\(b\in\Pi^{(2)}\), the distinguished real place contributes
\(\ell_x(b)/2\) to \(\cJ_{k_x}(\Tr_xb)\) by
\eqref{eq:distinguished-contribution}, while every local summand is
nonnegative.  The upper bound \eqref{eq:main-local} therefore forces
equality and makes every other local summand zero.  At every finite place,
\(\chi_v(\Tr_xb)=0\).  If
\(\lambda+\lambda^{-1}=\Tr_xb\), equation
\eqref{eq:lambda-chi} gives
\(\max\{|\lambda|_v,|\lambda|_v^{-1}\}=1\), so
\(|\lambda|_v=1\).  The ultrametric inequality yields
\(|\Tr_xb|_v=|\lambda+\lambda^{-1}|_v\le1\).
Thus all traces of the chosen lift \(\Delta_x\) are algebraic integers.
The full inverse image \(\widetilde\Lambda\) consists of the two lifts
\(\pm\rho_x(b)\), with \(b\in\Pi^{(2)}\).  Its traces are consequently
\(\pm\Tr_x b\); algebraic integrality is unchanged by sign, and the central
traces \(\pm2\) satisfy the same conclusion.

At every nondistinguished archimedean place, the roots in
\eqref{eq:z-equation} have modulus one.  Thus
\[
 \sigma(\Tr_xb)^2-2=z+z^{-1}\in[-2,2],
\]
so \(\sigma(\Tr_xb)^2\in[0,4]\).  It follows that
\(\sigma(\Tr_xb)\) is real and belongs to \([-2,2]\).  The same interval
contains its negative, so this verifies the conjugate-trace condition for
every element of the full inverse image \(\widetilde\Lambda\).  These traces
generate \(k_x\); hence every embedding of \(k_x\) has real image and
\(k_x\) is totally real.  All hypotheses of
Theorem~\ref{input:takeuchi} are now verified for the finite-covolume
Fuchsian group \(\Lambda\).  Takeuchi's theorem therefore implies that
\(\Lambda\) is arithmetic.  Arithmeticity is invariant under
commensurability, so \(\Gamma\) is arithmetic.

Suppose now that \(x\) is transcendental.  Then \(\dim Z_x>0\), and
Proposition~\ref{prop:bounded-degree-specialization} gives algebraic
Fuchsian characters \(y_j\in Z_x\) converging to \(x\), all in the same
Fuchsian component as \(x\).  Choose one of them.  Since \(x\) is
transcendental, \(y_j\ne x\).  Corollary~\ref{cor:specialization-rigidity}
therefore gives \(y_j=x\), a contradiction.  The transcendental case is therefore
impossible, and the algebraic argument proves the theorem.
\end{proof}

\begin{proof}[Proof of Corollary~\ref{cor:schmutz-sarnak}]
First suppose that \(\Gamma\) is torsion-free.  The critical-growth assertion
is Theorem~\ref{thm:main-arithmeticity}.  Linear growth implies the critical
bound.  If the signed trace set has bounded clustering, cover \([-X,X]\) by
\(O(X)\) intervals of radius one.  Each interval contains a uniformly
bounded number of trace values, so \eqref{eq:linear-growth} holds.

For a general cocompact Fuchsian lattice, Selberg's lemma supplies a
torsion-free finite-index subgroup \(\Gamma_0<\Gamma\)
\cite{AlperinSelberg}.  Its signed trace set is a subset of that of
\(\Gamma\), so each of the three hypotheses under consideration passes to
\(\Gamma_0\).  The torsion-free case makes \(\Gamma_0\) arithmetic.
Arithmeticity is a commensurability property, and therefore \(\Gamma\) is
arithmetic as well.
\end{proof}

\section{Dependencies among the estimates}
\label{sec:dependency-summary}

The following diagram records how the estimates are used in the proof.
Each arrow denotes a logical dependence.
\begin{small}
\[
\begin{gathered}
 \text{simultaneous-place Roth theorem}
 \longrightarrow
 \text{Corollary~\ref{cor:roth-recurrence-pair}},\\[2mm]
 \text{prime-square mixing and Lang--Weil}
 \longrightarrow
 \text{Lemma~\ref{lem:moving-bounded-exponents}},\\[2mm]
 \text{number-field divisor count}
 \longrightarrow
 \text{Lemma~\ref{lem:moving-two-trace-fibers}},\\[2mm]
 \text{bounded-exponent congruence counting}+
 \text{Lemma~\ref{lem:congruence-wedge}}+
 \text{Lemma~\ref{lem:moving-two-trace-fibers}}
 \longrightarrow
 \text{Theorem~\ref{thm:recurrence-divisor-upper}},\\[2mm]
 \text{Corollary~\ref{cor:roth-recurrence-pair}}+
 \text{Lemma~\ref{lem:moving-bounded-exponents}}+
 \text{Theorem~\ref{thm:recurrence-divisor-upper}}
 \longrightarrow
 \text{Theorem~\ref{thm:local-expansion}},\\[2mm]
 \text{algebraic specialization}+
 \text{Theorem~\ref{thm:local-expansion}}+
 \text{Thurston's asymmetric metric}
 \longrightarrow
 \text{\(x\) is algebraic},\\[2mm]
 \text{Theorem~\ref{thm:local-expansion} at \(x\)}+
 \text{local nonnegativity}
 \longrightarrow
 \text{trace integrality and bounded conjugates},\\
 \substack{
   \text{trace integrality and bounded conjugates}+
   \text{Takeuchi's criterion}\\
   {}+\text{commensurability invariance of arithmeticity}}
 \longrightarrow
 \text{Theorem~\ref{thm:main-arithmeticity}},\\[2mm]
 \text{Selberg's lemma}+\text{commensurability}
 \longrightarrow
 \text{Corollary~\ref{cor:schmutz-sarnak} for orbifolds}.
\end{gathered}
\]
\end{small}
The first five lines establish the local expansion inequality.  Algebraic
specialization then shows that the original character is algebraic, and
Takeuchi's criterion converts the local equality at that character into
arithmeticity.

\section*{Acknowledgements}

We thank Igor Rivin for helpful comments and a careful critical reading.

\bibliographystyle{siam}
\bibliography{biblio}

\end{document}
